\documentclass[10pt,a4paper]{amsart}
\usepackage[margin=19mm]{geometry}
\usepackage{amsmath,amssymb,mathtools}
\usepackage[scr=rsfs,cal=euler]{mathalfa}
\usepackage{xfrac}
\usepackage[unicode,hidelinks,bookmarks=false]{hyperref}
\newcommand{\ie}{i.e.\ }
\newcommand{\cf}{cf.\ }
\newcommand{\resp}{respectively\ }
\newcommand{\Subsection}{\subsection}
\providecommand{\nobreakdash}{\nobreak\hbox{-}}
\setlength{\emergencystretch}{2em}
\multlinegap0pt
\newtheorem{theorem}{Theorem}[section]
\newtheorem{lemma}[theorem]{Lemma}
\newtheorem{prop}[theorem]{Proposition}
\newtheorem{coro}[theorem]{Corollary}
\theoremstyle{definition}
\newtheorem{defi}[theorem]{Definition}
\newtheorem{rema}[theorem]{Remark}
\newtheorem{remas}[theorem]{Remarks}
\newtheorem{exam}[theorem]{Example}
\DeclareMathOperator{\Mult}{Mult}
\DeclareMathOperator{\irr}{irr}
\DeclareMathOperator{\lm}{lm}
\DeclareMathOperator{\Imrm}{Im}
\DeclareMathOperator{\card}{card}
%\DeclareMathOperator{\det}{det}
\DeclareMathOperator{\vect}{vect}
\DeclareMathOperator{\wt}{wt}


\newcommand*{\mk}{\mkern -1mu}
\newcommand*{\Mk}{\mkern -2mu}
\newcommand*{\mK}{\mkern 1mu}
\newcommand*{\MK}{\mkern 2mu}

\hypersetup{urlcolor=purple, linkcolor=blue, citecolor=red}


\newcommand*{\romanenumi}{\renewcommand*{\theenumi}{\roman{enumi}}}
\newcommand*{\Romanenumi}{\renewcommand*{\theenumi}{\Roman{enumi}}}
\newcommand*{\alphenumi}{\renewcommand*{\theenumi}{\alph{enumi}}}
\newcommand*{\Alphenumi}{\renewcommand*{\theenumi}{\Alph{enumi}}}

\title[Nonnegative k-Schur morphisms]{Alcove random walks and k-Schur functions: original Theorem5.6}
\author{Cédric Lecouvey and Pierre Tarrago}
\date{Source: Algebraic Combinatorics4(2)(2021),241--272; DOI10.5802/alco.147}
\begin{document}\maketitle
\noindent\textit{Editorial scope.} The counted unit is original Theorem5.6. Original sections2--5 supply its complete required k-Schur, restricted-graph, field-extension, Perron-Frobenius and boundary-extension core; all other results here remain uncounted. Later Markov-chain applications and further consequences are omitted. Ancillary links to Remark7.11 and Section6.3 point to the original article. Original mathematical expressions and hypotheses are retained without correction.
\setcounter{section}{1}
\section{Harmonic functions on the lattice of \texorpdfstring{$k$}{k}-bounded partitions}\label{section2}

\subsection{The lattices \texorpdfstring{$\mathcal{C}_{l}$}{Cl} and \texorpdfstring{$\mathcal{B}_{k}$}{BK}}
\label{subset_Lattices}

In this section, we refer to~\cite{LLMSSZ} and
\cite{Mac} for the material which is not defined. Fix $l>1$ a nonnegative
integer and set $k=l-1$. Let $\widetilde{W}$ be the affine Weyl group of type
$A_{k}^{(1)}$. As a Coxeter group, $\widetilde{W}$ is generated by the
reflections $s_{0},s_{1},\ldots,s_{k}$ so that its subgroup generated by
$s_{1},\ldots,s_{k}$ is isomorphic to the symmetric group $S_{l}$. Write
$\ell$ for the length function on $\widetilde{W}$. The group $\widetilde{W}$
determines a Coxeter arrangement by considering the hyperplanes orthogonal to
the roots of type $A_{k}^{(1)}$. The connected components of this hyperplane
arrangement yield a tessellation of $\mathbb{R}^{k}$ by alcoves on which the
action of $\widetilde{W}$ is regular. We denote by $A^{(0)}$ the fundamental
alcove. Write $\widetilde{R}$ for the set of affine roots of type $A_{k}%
^{(1)}$ and $R$ for its subset of classical roots of type $A_{k}$. The simple
roots are denoted by $\alpha_{0},\ldots,\alpha_{k}$ and $P$ is the weight
lattice of type $A_{k}$ with fundamental weights $\Lambda_{1},\ldots
,\Lambda_{k}$.

A reduced alcove path is a sequence of alcoves $(A_{1},\ldots,A_{m})$ such
that $A_{1}=A^{(0)}$ and for any $i=1,\ldots,m-1$, the alcoves $A_{i+1}$ and
$A_{i}$ share a common face contained in a hyperplane $H_{i}$ so that the
sequence $H_{1},\ldots,H_{m-1}$ is without repetition (each hyperplane can be
crossed only once). In the sequel, all the alcove paths we shall consider
will be reduced. For any $i=1,\ldots,m-1$, let $w_{i}$ be the unique element
of $\widetilde{W}$ such that $A_{i}=w_{i}(A^{(0)})$. Write $\vartriangleleft$
for the weak Bruhat order on $\widetilde{W}$ and $\rightarrow$ for the
covering relation $w\rightarrow w^{\prime}$ if and only if $w\vartriangleleft
w^{\prime}$ and $\ell(w^{\prime})=\ell(w)+1$. We then have $w_{1}\rightarrow
w_{2}\rightarrow\cdots\rightarrow w_{m}$.

\looseness-1
We shall identify a partition and its Young diagram. Recall that an $l$-core
can be seen as a partition where no box has hook length equal to $l$. Given an
$l$-core $\lambda$, we denote by $\ell(\lambda)$ its length which is equal to
the number of boxes of $\lambda$ with hook length less than $l$. Recall that
the residue of a box in a Young diagram is the difference modulo $l$ between
its row and column indices. We can define an arrow $\lambda\overset
{i}{\rightarrow}\mu$ between the two $l$-cores $\lambda$ and $\mu$ when
$\lambda\subset\mu$ and all the boxes in $\mu/\lambda$ have the same residue
$i$. By forgetting the label on the arrows, we get the structure of a graded rooted
graph $\mathcal{C}_{l}$ on the $l$-cores. For any two vertices $\lambda
\rightarrow\mu$ we have $\ell(\mu)=\ell(\lambda)+1$. Nevertheless, the
difference between the rank of the partitions $\lambda$ and $\mu$ is not
immediate to get in general.

The affine Grassmannian elements are the elements $w\in\widetilde{W}$ whose
associated alcoves are exactly those located in the fundamental Weyl chamber
(that is, in the Weyl chamber containing the fundamental alcove $A^{(0)}%
$). The $l$-cores are known to parametrize the affine Grassmannian elements.
More precisely, given two $l$-cores such that $\lambda\overset{i}{\rightarrow
}\mu$ and $w$ the affine Grassmannian element associated to $\lambda$,
$w^{\prime}=ws_{i}$ is the affine Grassmannian element associated to $\mu
$. In particular, we get $\ell(\lambda)=\ell(w)$. So reduced alcove paths in
the fundamental Weyl chamber, saturated chains of affine Grassmannian elements
and paths in $\mathcal{C}_{l}$ naturally correspond.

%\bigskip

A $k$-bounded partition is a partition $\lambda$ such that $\lambda_{1}\leq
k$. There is a simple bijection between the $l$-cores and the $k$-bounded
partitions. Start with a $l$ core $\lambda$ and delete all the boxes in the
diagram of $\lambda$ having a hook length greater than $l$ (recall there is no
box with hook length equal to $l$ since $\lambda$ is an $l$-core). This gives a
skew shape and to obtain a partition, move each row so obtained to the left.
The result is a $k$-bounded partition denoted $\mathfrak{p}(\lambda)$. For
some examples and the inverse bijection $\mathfrak{c}$, see~\cite[pages~18 and~19]{LLMSSZ}. This bijection permits to define an analogue of conjugation
for the $k$-bounded partitions. Given a $k$-bounded partition $\kappa$ set%
\[
\kappa^{\omega_{k}}=\mathfrak{p}(\mathfrak{c}(\kappa)^{\prime}).
\]
The graph $\mathcal{B}_{k}$ is the image of the graph $\mathcal{C}_{l}$ under
the bijection $\mathfrak{p}$. This means that $\mathcal{B}_{k}$ is the graph
obtained from $\mathcal{C}_{l}$ by deleting all the boxes with hook length
greater than~$l$ and next by aligning the rows obtained to the left. In
particular, reduced alcove paths in the fundamental Weyl chamber correspond to
$k$-bounded partitions paths in $\mathcal{B}_{k}$. We have the following lemma:

\begin{lemma}
We have an arrow $\kappa\rightarrow\delta$ in $\mathcal{B}_{k}$ if and only if
$\left\vert \delta\right\vert =\left\vert \kappa\right\vert +1,$
$\kappa\subset\delta$ and $\kappa^{\omega_{k}}\subset\delta^{\omega_{k}}%
$.\footnote{So $\mathcal{B}_{k}$ \emph{should not be confused} with the
subgraph of the Young graph with vertices the $k$-bounded partitions.}
\end{lemma}

Let $\Lambda$ be the algebra of symmetric functions in infinitely many
variables over $\mathbb{R}$. It is endowed with a scalar product $\langle
\cdot,\cdot\rangle$ such that $\langle s_{\lambda},s_{\mu}\rangle
=\delta_{\lambda,\mu}$ for any partitions $\lambda$ and $\mu$. Let
$\Lambda_{(k)}$ be the subalgebra of $\Lambda$ generated by the complete
homogeneous functions $h_{1},\ldots,h_{k}$. In particular, $\{h_{\lambda}%
\mid\lambda$ is $k$-bounded$\}$ is a basis of $\Lambda_{(k)}$.

\subsection{The \texorpdfstring{$k$}{k}-Schur functions}

We now define a distinguished basis of $\Lambda_{(k)}$ related to the graph
structures of $\mathcal{C}_{l}$ and $\mathcal{B}_{k}$. Consider $\lambda$ and
$\mu$ two $k$-bounded partitions with $\lambda\subset\mu$ and $r\leq k$ a
positive integer.

\begin{defi}
\label{Def_Hori_Strip}We will say that $\mu/\lambda$ is a weak horizontal
strip of size $r$ when

\begin{enumerate}
\item $\mu/\lambda$ is an horizontal strip with $r$ boxes (\ie the boxes in
$\mu/\lambda$ belong to different columns),

\item $\mu^{\omega_{k}}/\lambda^{\omega_{k}}$ is a vertical strip with $r$
boxes (\ie the boxes in $\mu^{\omega_{k}}/\lambda^{\omega_{k}}$ belong to
different rows).
\end{enumerate}
\end{defi}

Let us now define the notion of $k$-bounded semistandard tableau of shape
$\lambda$ a $k$-bounded partition and weight $\alpha=(\alpha_{1},\ldots
,\alpha_{d})$ a composition of $\left\vert \lambda\right\vert $ \emph{with no
part larger than }$k$. 

\begin{defi}
A $k$-bounded semistandard tableau of shape $\lambda$ is a semistandard
filling of $\lambda$ with integers in $\{1,\ldots,d\}$ such that for any
$i=1,\ldots,d$ the boxes containing $i$ define a weak horizontal strip of size
$\alpha_{i}$.
\end{defi}

One can prove that for any $k$-bounded partitions $\lambda$ and $\alpha$ the
number $K_{\lambda,\alpha}^{(k)}$ of $k$-bounded semistandard tableaux of
shape $\lambda$ and weight $\alpha$ verifies%
\[
K_{\lambda,\lambda}^{(k)}=1\text{ and }K_{\lambda,\alpha}^{(k)}\neq
0\Longrightarrow\alpha\leq\lambda
\]
where $\leq$ is the dominant order on partitions.

\begin{defi}
The $k$-Schur functions $s_{\kappa}^{(k)},\kappa\in\mathcal{B}_{k}$ are the
unique functions in $\Lambda_{(k)}$ such that
\[
h_{\delta}=\sum_{\delta\leq\kappa,\kappa\in\mathcal{B}_{k}}K_{\kappa,\delta
}^{(k)}s_{\kappa}^{(k)}%
\]
for any $\delta$ in $\mathcal{B}_{k}$.
\end{defi}

\begin{prop}[Pieri rule for $k$-Schur functions]
For any $r\leq k$ and any $\kappa
\in\mathcal{B}_{k}$ we have
\begin{equation}
h_{r}s_{\kappa}^{(k)}=\sum_{\varkappa\in\mathcal{B}_{k}}s_{\varkappa}^{(k)}
\label{PieriKSchur}%
\end{equation}
where the sum is over all the $k$-bounded partitions $\varkappa$ such that
$\varkappa/\kappa$ is a weak horizontal strip of size $r$ in $\mathcal{B}_{k}$.
\end{prop}

When $r=1$, the multiplication by $h_{1}$ is easily described by considering
all the possible $k$-bounded partitions at distance $1$ from $\kappa$ in
$\mathcal{B}_{k}$. Thanks to a geometric interpretation of the $k$-Schur
functions in terms of the homology of affine Grassmannians, Lam showed that
the product of two $k$-Schur functions is $k$-Schur positive:

\begin{theorem}[{\cite[Corollary~8.2]{LamLR}}] Given $\kappa$ and $\delta$ two $k$-bounded partitions, we have
\[
s_{\kappa}^{(k)}s_{\delta}^{(k)}=\sum_{\nu\in\mathcal{B}_{k}}c_{\lambda
,\delta}^{\nu(k)}s_{\nu}^{(k)}%
\]
with $c_{\lambda,\delta}^{\nu(k)}\in\mathbb{Z}_{\geq0}$.
\end{theorem}

\subsection{Some properties of \texorpdfstring{$k$}{K}-Schur functions}

The $k$-conjugation operation $\omega_{k}$ can be read directly at the level
of $k$-bounded partitions without using the ordinary conjugation operations on
the $l$-cores (see (1.9)) in~\cite{LLMSSZ}). To do this, start with a
$k$-bounded partition $\lambda=(\lambda_{1},\ldots,\lambda_{r})$ and decompose
it into its chains $\{c_{1},c_{2},\ldots,c_{r}\}$ where each chain is a
sequence of parts of $\lambda$ obtained recursively as follows. The procedure
is such that any part $\lambda_{i}$ is in the same chain as the part
$\lambda_{i+k-\lambda_{i}+1}$ when $i+k-\lambda_{i}+1\leq r$ (from the part
$\lambda_{i}$ one jumps $k-\lambda_{i}$ parts to get the following part of the
chain). Observe in particular that all the parts with length $k$ belong to the
same chain for in this case we jump $0$ parts. Once the chains $c_{i}$ are
determined, $\lambda^{\omega_{k}}$ is the partition with $k$-columns whose
lengths are the sums of the $c_{i}$'s$.$

\begin{exam}
Consider the $5$-partition $\lambda=(\boldsymbol{5,5,5,4},$\emph{4}%
$,\boldsymbol{3},$\emph{3}$,3,\boldsymbol{2},$\emph{2}$,1)$. Then, we get
$c_{1}=\{5,5,5,4,3,2\}$ next $c_{2}=\{4,3,2\},$ $c_{3}=\{3,1\},$
$c_{4}=\emptyset$ and $c_{5}=\emptyset$. So $\lambda^{\omega_{5}}$ is the
partition with columns of heights $24,$ $9$ and $4$.
\end{exam}

The following facts will be useful.

\begin{enumerate}
\item Any partition $\lambda$ of rank at most $k$ is a $k$-bounded partition
and is then equal to its associated $l$-core (because $\lambda$ has no hook of
length $l=k+1$).

\item The lattice $\mathcal{B}_{k}$ coincides with the ordinary Young lattice
on the partitions of rank at most $k$. On this subset $\omega_{k}$ is the
ordinary conjugation.

\item For any partition $\lambda$ of rank at most $k,$ the $k$-Schur function
coincides with the ordinary Schur function that is $s_{\lambda}^{(k)}%
=s_{\lambda}$. In particular, the homogeneous functions $h_{1},\ldots,h_{k}$
and the elementary functions $e_{1},\ldots,e_{k}$ are the $k$-Schur functions
corresponding to the rows and columns partitions with at most $k$ boxes, respectively.
\end{enumerate}

The $k=2$ case is easily tractable because the lattice of $2$-bounded
partitions we consider has a simple structure. One verifies easily that for
any $2$-bounded partition $\lambda=(2^{a},1^{n-2a})$, we get in that case
\[
s_{\lambda}^{(2)}=
\begin{cases}
h_{2}^{a}e_{2}^{\frac{n}{2}-a}&\text{ when }n\text{ even,}\\[4pt]
h_{2}^{a}e_{2}^{\frac{n-1}{2}-a}e_{1}&\text{ when }n\text{ is odd.}%
\end{cases}
\]
When $k>2,$ the structure of the graph $\mathcal{B}_{k}$ becomes more
complicated. Given a $k$-bounded partition $\lambda$ one can first precise
where it is possible to add a box in the Young diagram of $\lambda$ to get an
arrow in $\mathcal{B}_{k}$. Assume we add a box in the row $\lambda_{i}$ of
$\lambda$ to get the $k$-bounded partition $\mu$, denote by $c=\{\lambda
_{i_{1}},\ldots,\lambda_{i_{r}}\}$ the chain containing $\lambda_{i}$ where we
have $\lambda_{i}=\lambda_{i_{a}}$ with $a\in\{1,\ldots,r\}$. Observe we can
add components equal to zero to $c$ if needed since $\lambda$ is defined up to
an arbitrary number of zero parts. The following lemma permits to avoid the
use of $\omega_{k}$ in the construction of $\mathcal{B}_{k}$.

\begin{lemma}
\label{conditionArrow} There is an arrow $\lambda\rightarrow\mu$ in
$\mathcal{B}_{k}$ if and only if $\lambda_{i_{b}-1}=\lambda_{i_{b}}$ for any
$b=a+1,\ldots,r,$ that is if and only if each part located up to $\lambda_{i}$
in the chain containing $\lambda_{i}$ is preceded by a part with the same size.
\end{lemma}

\begin{proof}
One verifies that if the previous condition is not satisfied, $\mu^{\omega
_{k}}$ and $\lambda^{\omega_{k}}$ will differ by at least two boxes and if it
is satisfied by only one as desired.
\end{proof}

%\bigskip

\begin{exam}\ \\*[-1.2em]

\begin{enumerate}
\item One can always add a box in the first column of $\lambda$ since the
parts located up to the part $0$ are all equal to $0$.

\item Assume $k=2$, then we can add a box to the part $\lambda_{i}$ equal to
$1$ to get a part equal to $2$ if and only if there is an even number of parts
equal to $1$ up to $\lambda_{i}$. This is equivalent to saying that $\lambda$
has an odd number of parts equal to $1$, that is that the rank of $\lambda$ is
odd since the other parts are equal to $2$ (or $0$).
\end{enumerate}
\end{exam}

%\bigskip

For any $a=1,\ldots,k$, let $R_{a}$ be the rectangle partition $(k-a+1)^{a}$.
The previous observations can be generalized (see~\cite[Corollary~8.3]{LamLR}  and~\cite[Section~4]{LLMSSZ}):

\begin{prop}\ \\*[-1.2em] 

\begin{enumerate}
\item Assume $\lambda$ is a $k$-bounded partition which is also a
$(k+1)$-core. Then $s_{\lambda}^{(k)}=s_{\lambda}$ (that is the $k$-Schur and
the Schur functions corresponding to $\lambda$ coincide).

\item In particular, for any rectangle partition $R_{a}$, we have $s_{R_{a}%
}^{(k)}=s_{R_{a}}$.

\item For any $a=1,\ldots,k$ and any $k$-bounded partition $\lambda$ we have
\[
s_{R_{a}}s_{\lambda}^{(k)}=s_{\lambda\cup R_{a}}^{(k)}%
\]
where $\lambda\cup R_{a}$ is obtained by adding $a$ times a part $k-a+1$ to
$\lambda$.
\end{enumerate}
\end{prop}

\begin{coro}\label{reducedCase} 
\looseness-1
For each $k$-bounded partition $\lambda$, there exists a
unique irreducible $k$-bounded partition\footnote{A $k$-bounded partition is
called irreducible when it contains less or equal to $a$ parts equal to $k-a$ for any
$a=0,\ldots,k-1$. Otherwise it is called reducible.} $\widetilde{\lambda}$ and a unique sequence of nonnegative
integers $p_{1},\ldots,p_{k}$ such that
\[
s_{\lambda}^{(k)}=\prod_{a=1}^{k}s_{R_{a}}^{p_{a}}s_{\widetilde{\lambda}%
}^{(k)}.
\]
In particular, the $k$-Schur functions are completely determined by the
$k$-Schur functions indexed by the irreducible $k$-bounded partitions and by
the ordinary Schur functions $s_{(k-a+1)^{a}}$, $a=1,\ldots,k$.
\end{coro}

\begin{rema}
\label{Rem_bij}Write $\mathcal{P}_{\irr }$ for the set of irreducible
partitions. The map $\Delta:$ $\mathcal{B}_{k}\rightarrow\mathcal{P}%
_{\irr }\times\mathbb{Z}_{\geq0}^{k}$ which associates to each
$k$-bounded partition $\lambda$ the pair $(p,\widetilde{\lambda})$ where
$p=(p_{1},\ldots,p_{k})$ is a bijection.
\end{rema}

\begin{exam}
Assume $k=4$ and $\lambda=(4,4,3,3,3,3,3,2,2,2,2,1,1,1,1,1)$. Then we get
$\widetilde{\lambda}=(3,2,1)$ and
\[
s_{\lambda}^{(k)}=s_{(4)}^{2}s_{(3,3)}^{2}s_{(2,2,2)}s_{(1,1,1,1)}%
s_{\widetilde{\lambda}}^{(k)}.
\]

\end{exam}

We conclude this paragraph by recalling other important properties of
$k$-Schur functions. We have first the inclusions of algebras $\Lambda
_{(k)}\subset\Lambda_{(k+1)}\subset\Lambda$.

\begin{prop}
[{see~\cite[Section~4]{LLMSSZ}}] \label{Prop_positiveExp}\ \\*[-1.2em] 

\begin{enumerate}
\item Each $k$-Schur function has a positive expansion on the basis of
$(k+1)$-Schur functions.

\item Each $k$-Schur function has a positive expansion on the basis of
ordinary Schur functions.
\end{enumerate}
\end{prop}

\subsection{Harmonic functions and minimal boundary of \texorpdfstring{$\mathcal{B}_{k}$}{Bk}}

\begin{defi}
A function $f:\mathcal{B}_{k}\rightarrow\mathbb{R}$ is said harmonic when
\[
f(\lambda)=\sum_{\lambda\rightarrow\mu}f(\mu)\text{ for any }\lambda
\in\mathcal{B}_{k}.
\]
We denote by $\mathcal{H}(\mathcal{B}_{k})$ the set of harmonic functions on
$\mathcal{B}_{k}$.
\end{defi}

Another way to understand harmonic functions is to introduce the infinite
matrix $\mathcal{M}$ of the graph $\mathcal{B}_{k}$. The harmonic functions on
$\mathcal{B}_{k}$ then correspond to the right eigenvectors for $\mathcal{M}$
associated to the eigenvalue $1$. One can also consider $t$-harmonic functions
which correspond to the right eigenvectors for $\mathcal{M}$ associated to the
eigenvalue $t$. Clearly $\mathcal{H}(\mathcal{B}_{k})$ is a vector space over
$\mathbb{R}$. In fact, we mostly restrict ourself to the set $\mathcal{H}%
^{+}(\mathcal{B}_{k})$ of positive harmonic functions for which $f$ takes
values in $\mathbb{R}_{\geq0}$. Then, $\mathcal{H}^{+}(\mathcal{B}_{k})$ is a
cone since it is stable under addition and multiplication by a positive real. To
study $\mathcal{H}^{+}(\mathcal{B}_{k})$, we only have to consider its subset
$\mathcal{H}_{1}^{+}(\mathcal{B}_{k})$ of normalized harmonic functions such
that $f(1)=1$. In fact, $\mathcal{H}_{1}^{+}(\mathcal{B}_{k})$ is a convex set
and its structure is controlled by its extremal subset $\partial
\mathcal{H}^{+}(\mathcal{B}_{k})$.

We aim to characterize the extremal positive harmonic functions
defined on $\mathcal{B}_{k}$ and obtain a simple parametrization of
$\partial\mathcal{H}^{+}(\mathcal{B}_{k})$. By using the Pieri rule on
$k$-Schur functions, we get%
\[
s_{\lambda}s_{(1)}=\sum_{\lambda\rightarrow\mu}s_{\mu}%
\]
for any $k$-bounded partitions $\lambda$ and $\mu$. This means that
$\mathcal{B}_{k}$ is a so-called multiplicative graph with associated algebra
$\Lambda_{(k)}$. Moreover, if we denote by $K$ the positive cone spanned by
the set of $k$-Schur functions, we can apply the ring theorem of Kerov and
Vershik (see for example~\cite[Section~8.4]{LT}) which characterizes the
subset of extreme points $\partial\mathcal{H}^{+}(\mathcal{B}_{k})$. Denote
by $\Mult ^{+}(\Lambda_{(k)})\subset\Lambda_{(k)}^{\ast}$ the set of
multiplicative functions on $\Lambda_{(k)}$ which are nonnegative on $K$ and
equal to $1$ on $s_{1}$. So a function $f$ belongs to $\Mult %
^{+}(\Lambda_{(k)})$ when $f:\Lambda_{(k)}\rightarrow\mathbb{R}$ is linear and
satisfies $f(UV)=f(U)f(V)$ for any $U,V\in\Lambda_{(k)}$. Note that
$i:\mathcal{B}_{k}\longrightarrow\Lambda_{(k)}$ such that $i(\lambda
)=s_{\lambda}^{(k)}$ induces a map $i^{\ast}:\Lambda_{(k)}^{\ast
}\longrightarrow F(\mathcal{B}_{k},\mathbb{R})$. Let $K.K$ be the set of products of two elements in $K$. Since we have $K.K\subset K$,
we get the following algebraic characterization of $\partial\mathcal{H}%
^{+}(\mathcal{B}_{k})$.

\begin{prop}
\label{multiGraphExtreme}The map $i^{\ast}$ yields an homeomorphism between
$\Mult ^{+}(\Lambda_{(k)})$ and $\partial\mathcal{H}^{+}(\mathcal{B}%
_{k})$. 
\end{prop}

Since $i(\mathcal{B}_{k})$ is a basis of $\Lambda_{(k)}$, this means that
$\partial\mathcal{H}(\mathcal{B}_{k})$ is completely determined by the
$\mathbb{R}$-algebra morphisms $\varphi:\Lambda_{(k)}\rightarrow\mathbb{R}$
such that $\varphi(s_{1})=1$ and $\varphi(s_{\lambda}^{(k)})\geq0$ for any
$k$-bounded partition $\lambda$. Each function $f\in\partial\mathcal{H}%
^{+}(\mathcal{B}_{k})$ can then be written $f=\varphi\circ i$.

By Corollary~\ref{reducedCase}, the condition $\varphi(s_{\lambda}^{(k)}%
)\geq0$ for each $k$-bounded partition reduces in fact to test a finite number
of $k$-Schur functions, namely $\varphi(s_{\widetilde{\lambda}}^{(k)})\geq0$
for each irreducible $k$-bounded partition (there are $k!$ such
partitions) and $\varphi(s_{(k-a+1)^{a}})\geq0$ for any $a=1,\ldots,k$. We
will see that the condition $\varphi(s_{\lambda}^{(k)})>0$ for each
$k$-bounded partition reduces in fact to test the positivity on the Schur
functions associated to partitions with maximal hook-length less or equal to
$k$ (See Remark~\href{https://alco.centre-mersenne.org/articles/10.5802/alco.147/}{7.11}).

\section{Restricted graph and irreducibility}\label{section3}

\subsection{The matrix \texorpdfstring{$\Phi$}{Phi}}

\label{subsec_MatrixFI}By Corollary~\ref{reducedCase}, each morphism
$\varphi:\Lambda_{(k)}\rightarrow\mathbb{R}$ is uniquely determined by its
values on the rectangle Schur functions $s_{R_{a}},1\leq a\leq k$ and on each
$s_{\tilde{\lambda}}^{(k)}$ where $\tilde{\lambda}$ is an irreducible
$k$-bounded partition. Set $r_{a}=\varphi(s_{R_{a}}),a=1,\ldots,k$ and
$\vec{r}=(r_{1},\ldots,r_{k})$. Recall that $\mathcal{P}_{\irr }$ is
the set of irreducible $k$-bounded partitions (including the empty partition).
Then, for $\lambda\in\mathcal{P}_{\irr }$,
\begin{equation}
\varphi(s_{\lambda}^{(k)})\varphi(s_{(1)})=\sum_{\lambda\rightarrow\mu}%
\varphi(s_{\mu}^{(k)}). \label{PieriPhi}%
\end{equation}
By Corollary~\ref{reducedCase}, for each $k$-bounded partition $\mu$ there
exists a sequence $\{p_{1}^{\mu},p_{2}^{\mu}\ldots,p_{k}^{\mu}\}$ of elements
in $\{1,\dots,k\}$ and an irreducible partition $\widetilde{\mu}$ such that
\begin{equation}
s_{\mu}^{(k)}=\prod_{a=1}^{k}s_{R_{a}}^{p_{a}^{\mu}}s_{\widetilde{\mu}}%
^{(k)}\text{ and thus }\varphi(s_{\mu}^{(k)})=\prod_{a=1}^{k}r_{a}^{p_{a}%
^{\mu}}\varphi(s_{\widetilde{\mu}}^{(k)}). \label{decompos}%
\end{equation}
Hence by setting
\[
\varphi_{\lambda\nu}=\sum_{\substack{\lambda\rightarrow\mu\\\widetilde{\mu
}=\nu}}\prod_{1\leq a\leq k}r_{a}^{p_{a}^{\mu}}%
\]
we get
\[
\varphi(s_{\lambda}^{(k)})=\sum_{\nu\in\mathcal{P}_{\irr }}%
\varphi_{\lambda\nu}\varphi(s_{\nu}^{(k)}).
\]
Let $\Phi_{(r_{1},\ldots,r_{k})}:=(\varphi_{\nu\lambda})_{\lambda,\nu
\in\mathcal{P}_{\irr }}$ \footnote{Observe we have defined
$\Phi_{(r_{1},\ldots,r_{k})}$ as the transpose of the matrix $(\varphi
_{\lambda,\mu})_{\lambda,\nu\in\mathcal{P}_{\irr }}$ to make it
compatible with the multiplication by $s_{(1)}$ used in Section
\ref{Secion_Field}.} and define $f\in\mathbb{R}^{\mathcal{P}_{\irr }}$
as the vector $(\varphi(s_{\lambda}^{(k)}))_{\lambda\in\mathcal{P}%
_{\irr }}$. When there is no risk of confusion, we simply write $\Phi$
instead of $\Phi_{(r_{1},\ldots,r_{k})}$. The vector $f$ is a left eigenvector
of $\Phi$ for the eigenvalue $\varphi(s_{1})$ with positive entries having
value $1$ on $\emptyset$ and $\varphi(s_{1})$ on $s_{1}$.

\subsection{Irreducibility of the matrix \texorpdfstring{$\Phi$}{Phi}}

Recall that a matrix $M\in M_{n}(\mathbb{R})$ with nonnegative entries is
irreducible if and only if for each $1\leq i,j\leq n$ there exists $n\geq1$
such that $(M^{n})_{ij}>0$.

\begin{prop}
\label{irredPhi}Consider $\vec{r}=(r_{1},\ldots,r_{k})\in\mathbb{R}_{\geq
0}^{k}.$ Then, the matrices $\Phi_{(r_{1},\ldots,r_{k})}$ and $\Phi
_{(r_{1},\ldots,r_{k})}^{t}$ associated to $\varphi$ are irreducible if and
only if for all $1\leq a\leq k-1$, $r_{a}$ or $r_{a+1}$ is positive.
\end{prop}

We will prove in fact that $\Phi^{t}$ is irreducible. Let $G$ be the graph
with set of vertices $\mathcal{P}_{\irr }$ and a directed edge from
$\lambda$ to $\nu$ if and only if $\Phi_{\lambda\nu}\not =0$. The matrix
$\Phi^{t}$ is irreducible if and only if $G$ is strongly connected, which
means that there is a (directed) path from any vertex to any other vertex of
the graph. We prove Proposition~\ref{irredPhi} by showing that $G$ is strongly
connected. Let us first establish a preliminary lemma. We say that $\lambda
\in\mathcal{P}_{\irr }$ is $1$-saturated when it contains less than
$k-1$ parts equal to $1$. More generally, for $i\geq2$ the partition
$\lambda\in\mathcal{P}_{\irr }$ is $i$-saturated when we have%
\[
\lambda=(\ldots,(i-1)^{k-i+1},\dots,2^{k-2},1^{k-1})\text{ and }\lambda
\neq(\ldots,i^{k-i},(i-1)^{k-i+1},\dots,2^{k-2},1^{k-1}).
\]
Denote by $\lambda^{1}$ the irreducible $k$-bounded partition $((k-1)^{1}%
,(k-2)^{2},\dots,1^{k-1})$. Observe that $\lambda\in\mathcal{P}_{\irr %
}$ is $k$-saturated if and only if $\lambda=\lambda^{1}$.

\begin{lemma}
Any vertex $\lambda$ of $G$ is connected to $\lambda^{1}$.
\end{lemma}

\begin{proof}
We get a path between $\lambda$ and $\lambda^{1}$ from the following observations.

\noindent1: If $\lambda$ is $i$-saturated, we claim that $\lambda
\rightarrow\lambda^{\uparrow i-1}$, where $\lambda^{\uparrow i-1}$ is the
partition obtained by adding one box to the first row of size $(i-1)$ of
$\lambda$. Moreover, $\lambda^{\uparrow i-1}$ is then $(i-1)$-saturated. To
see this, consider the minimal integer $r$ such that $\lambda_{r}=i-1$ and the
chain $c$ associated to $\lambda_{r}$. Since $\lambda$ is $i$-saturated, a
combinatorial computation shows that
\[
c=\{\ldots,\lambda_{r},\lambda_{r+k-i+2},\lambda_{r+(k-i+2)+(k-i+3)}%
,\ldots,\lambda_{m}\}\text{ with }m=r+\sum_{s=0}^{i-1}(k-i+s).
\]
Moreover, for $0\leq t\leq i-1$, the $(t+1)$-th row in the chain $c$ is
$\lambda_{r+\sum_{s=0}^{t}(k-i+s)}$ and has length $i-1-t$. In particular,
each row following $\lambda_{r}$ in the chain $c$ is preceded by a row with
the same length. Therefore, we can apply Lemma~\ref{conditionArrow} to get the
existence of an edge between $\lambda$ and the partition $\lambda^{\uparrow
i-1}$ which is $(i-1)$-saturated.

\noindent2: Suppose that $\lambda$ is $i$-saturated. Applying successively the
previous procedure on the partitions $\lambda,\lambda^{\uparrow(i-1)}%
,(\lambda^{\uparrow(i-1)})^{\uparrow(i-2)},\dots$ eventually yields a
partition $\nu$ which is irreducible, $i$-saturated, with one more row of
length $i$ than $\lambda$ and such that there is a path between $\lambda$ and
$\nu$ in $G$.

\noindent3: Suppose that $\lambda$ has initially $l\leq k-i-1$ rows of size
$i$. By repeating $k-i-l$ times step 2, one gets a partition $\kappa$
connected to $\lambda$ in $\Gamma$ which is $(i+1)$-saturated.

\noindent
4: Repeated applications of step 3 for $i<l\leq k$ yields a path
between $\lambda$ and $\lambda^{1}$ in~$G$.
\end{proof}

We prove Proposition~\ref{irredPhi} by giving necessary and sufficient
conditions to have a path in $G$ between $\lambda^{1}$ and $\emptyset$.

\begin{proof}
[Proof of Proposition~\ref{irredPhi}]Let us show that there is a path between
$\lambda^{1}$ and $\emptyset$ if and only if for all $1\leq a\leq k-1$,
$\varphi(s_{R_{a}})$ or $\varphi(s_{R_{a+1}})$ is positive.

Assume first that for all $1\leq a\leq k-1$, $\varphi(s_{R_{a}})$ or
$\varphi(s_{R_{a+1}})$ is positive.

\noindent For $1\leq a\leq k-1$, let $\lambda^{a}$ be the partition
$(k-1,(k-2)^{2},\ldots,a^{k-a})$, and set $\lambda^{k}=\emptyset$. Let us
prove that for $1\leq a\leq k-1$, there is a path in $G$ from $\lambda^{a}$ to
$\lambda^{a+1}$. If $\varphi(s_{R_{a}})>0$, it suffices to add a part of
length $a$ at the end of $\lambda^{a}$, which is always possible. This gives a
rectangle $a^{k-a+1}$ which, once removed, yields $\lambda^{a+1}$. Assume now
we have $\varphi(s_{R_{a}})=0$ and thus $\varphi(s_{R_{a+1}})>0$ by the
hypothesis on $\varphi$. Let $i$ be minimal such that $\lambda_{i}^{a}=a$, and
let $c$ be the chain containing $\lambda_{i}^{a}$. On the one hand, since
$\lambda_{i}^{a}=a$, the part following $\lambda_{i}^{a}$ in $c$ is
$\lambda_{i+k-a+1}^{a}$. On the other hand, by definition of $\lambda^{a}$, we
have $\lambda_{i+(k-a)}^{a}=\lambda_{i+(k-a)+1}^{a}=0$. Thus, by Lemma
\ref{conditionArrow}, we can add a box on the part $\lambda_{i}^{a}$, which
makes appear a block $(a+1)^{k-a+1}$. Since $\varphi(s_{R_{a+1}})>0$ we so get
an arrow in $G$ between $\lambda^{a}$ and the partition $(\lambda
^{a-2},\lambda_{i+1}^{a},\ldots,\lambda_{i+(k-a)-1}^{a})$. Similarly, we can
successively add a box to the parts $\lambda_{i+1}^{a},\ldots,\lambda
_{i+(k-a)-1}^{a}$ (all equal to $a$) and get a path between $(\lambda
^{a-2},\lambda_{i+1}^{a},\ldots,\lambda_{i+(k-a)-1}^{a})$ and $\lambda^{a+1}$.
We so obtain a path in $G$ from $\lambda^{a}$ to $\lambda^{a+1}$ for all
$1\leq a\leq k-1$ and thus a path from $\lambda^{1}$ to $\lambda^{k}%
=\emptyset$.

Assume now that there exists $1\leq a\leq k-1$ such that $\varphi(s_{R_{a}%
})=\varphi(s_{R_{a+1}})=0$.

\noindent
Let $\gamma=(\mu^{1},\ldots,\mu^{r})$ be a path in $G$ starting at
$\mu^{1}=a^{k-a}$, and denote by $x_{i}$ and $y_{i}$ the number of parts in
$\mu^{i}$ equal to $a+1$ and $a$, respectively. Since $\gamma$ is a path in
$G$, for all $1\leq i\leq r$, we have $x_{i}\leq k-a-1$ and $y_{i}\leq k-a$.
Let us prove by induction on $1\leq i\leq r$ that $x_{i}+y_{i}\geq k-a$.

\noindent This is certainly true for $i=1$. Assume that $i>1$ and the result
holds for $i-1$. Since $\varphi(s_{R_{a}})=\varphi(s_{R_{a+1}})=0$, the only
way to get $x_{i}+y_{i}<x_{i-1}+y_{i-1}$ is to add a box in the first part of
length $a+1$ (if any) of $\mu^{i-1}$. So if we have $x_{i-1}=0$, then
$x_{i}+y_{i}\geq x_{i-1}+y_{i-1}\geq k-a$. Assume now that $x_{i-1}>0$. Let
$l$ be minimal such that $\mu_{l}^{i-1}=a+1$ and let $c$ be the chain
containing $\mu_{l}^{i-1}$. The part following $\mu_{l}^{i-1}$ in $c$ is then
equal to $l+k-(a+1)+1=l+(k-a)$. Since $x_{i-1}+y_{i-1}\geq k-a$ and
$x_{i-1}\leq k-a-1$, we have $\mu_{l+(k-a)-1}^{i-1}=a$. Thus, by Lemma
\ref{conditionArrow}, it is possible to add a box on the part $l$ of
$\mu^{i-1}$ if and only if $\mu_{l+k-a}^{i-1}=a.$ If so, we get in fact
$x_{i-1}+y_{i-1}\geq(l+k-a)-l+1\geq k-a+1$ and $x_{i}+y_{i}=x_{i-1}+y_{i-1}%
-1$. Therefore, in any cases, $x_{i}+y_{i}\geq k-a$. We have proved that for
any path in $G$ starting at $a^{k-a}$, the number of parts of length $a$ or
$a+1$ is at least $k-a$, thus there is no path between $a^{k-a}$ and
$\emptyset$ in $G$.
\end{proof}

\section{Field extensions and \texorpdfstring{$k$}{k}-Schur functions}\label{section4}

\label{Secion_Field}

\subsection{Field extensions}

Recall that $\Lambda_{(k)}\mathbb{=R}[h_{1},\ldots,h_{k}]$. Since
$h_{1},\ldots,h_{k}$ are algebraically independent over $\mathbb{R}$, we can
consider the fraction field $\mathbb{L}=\mathbb{R}(h_{1},\ldots,h_{k}%
)$. Write $\mathbb{A}=\mathbb{R}[s_{R_{1}},\ldots s_{R_{k}}]$ the subalgebra
of $\Lambda_{(k)}$ generated by the rectangle Schur functions $s_{R_{a}%
},a=1,\ldots,k$. In order to introduce the fraction field of $\mathbb{A}$, we
first need to check that $s_{R_{1}},\ldots s_{R_{k}}$ are algebraically
independent over $\mathbb{R}$. We shall use a proposition giving a sufficient
condition on a family of polynomials to be algebraically independent.

Let $k$ be a field and $k[T_{1},\ldots,T_{m}]$ the ring of polynomials in
$T_{1},\ldots,T_{m}$ over $k$. For any $\beta\in\mathbb{Z}_{\geq0}^{m}$, we
set $T^{\beta}=T_{1}^{\beta_{1}}\cdots T_{m}^{\beta_{m}}$. We also consider a total order $\preceq$ on the monomials of $k[T_{1},\ldots,T_{m}]$ given by the lexicographical order on the exponents. Namely, $T^{\beta}\succeq T^{\beta'}$ if and only if $\beta_{i}> \beta'_{i}$ on the first index on which $\beta$ and $\beta'$ differ, if any. The
leading monomial $\lm (P)$ of a polynomial $P\in k[T_{1},\ldots,T_{m}]$
is the monomial appearing in the support of $P$ (that is, with a nonzero
coefficient) maximal under the total order $\preceq$.

\begin{prop}
[{See~\cite[Lemma~4.2.10]{Mitt},~\cite[Proposition~6.6.11]{KR}}]\label{Prop_alge_inde}\ \\*[-1.2em]

\begin{enumerate}
\item\label{prop4.1_1} The monomials $T^{\beta^{(1)}},\ldots,T^{\beta^{(p)}}$ are algebraically
independent if and only if $\beta^{(1)},\ldots,\beta^{(p)}$ are linearly
independent over $\mathbb{Z}$.

\item\label{prop4.1_2} Consider $P_{1},\ldots,P_{l}$ polynomials in $k[T_{1},\ldots,T_{m}]$
such that $\lm (P_{1}),\ldots,\lm (P_{l})$ are algebraically
independent, then $P_{1},\ldots,P_{l}$ are algebraically independent.
\end{enumerate}
\end{prop}
We also refer to~\cite[Section~3]{Pan} for a detailed proof of this proposition. With Proposition~\ref{Prop_alge_inde} in hand, it is then easy to check that
$s_{R_{1}},\ldots s_{R_{k}}$ are algebraically independent over $\mathbb{R}$.
Recall that each Schur function $s_{\lambda}$ with indeterminate set
$X=\{X_{1},X_{2}\ldots\}$ decomposes on the form%
\[
s_{\lambda}=X^{\lambda}+\sum_{\mu<\lambda}K_{\lambda,\mu}X^{\mu}%
\]
where $\leq$ is the dominant order over finite sequences of integers, that is,
$\beta<\beta^{\prime}$ when $\beta^{\prime}-\beta$ decomposes as a sum of
$\varepsilon_{i}-\varepsilon_{j},i<j$ with nonnegative integer coefficients.
Since the lexicographical order refines the dominance order,
by Assertion~\eqref{prop4.1_1} of the previous proposition, the monomials $X^{R_{1}}%
,\ldots,X^{R_{k}}$ are algebraically independent since the rectangle
partitions $R_{a},a=1,\ldots,k$ are linearly independent over $\mathbb{Z}$.
Assertion~\eqref{prop4.1_2} then implies that $s_{R_{1}},\ldots s_{R_{k}}$ are algebraically
independent over $\mathbb{R}$. We denote by $\mathbb{K}=\mathbb{R}(s_{R_{1}%
},\ldots s_{R_{k}})$ the fraction field of the algebra $\mathbb{A}$.

\begin{prop}\label{PropLK}\ \\*[-1.2em]

\begin{enumerate}
\item\label{prop4.2_1} Each $h_{1},\ldots,h_{k}$ is algebraic over $\mathbb{K}$.

\item\label{prop4.2_2} We have $\mathbb{L}=\mathbb{R}(h_{1},\ldots,h_{k})=\mathbb{K}%
[h_{1},\ldots,h_{k}]$.

\item\label{prop4.2_3} $\mathbb{L}$ is an algebraic extension of $\mathbb{K}$.

\item\label{prop4.2_4} The field $\mathbb{L}$ is a finite extension of $\mathbb{K}$ with degree
$[\mathbb{L}:\mathbb{K}]=k!$ and the set $\mathcal{I}=\{s_{\kappa}^{(k)}%
\mid\kappa\in\mathcal{P}_{\irr }\}$ is a basis of $\mathbb{L}$ over
$\mathbb{K}$.

\item\label{prop4.2_5} $\Lambda_{(k)}$ is an integral extension of $\mathbb{A}$.
\end{enumerate}
\end{prop}

\begin{proof}
\eqref{prop4.2_1}:~For any $a=1,\ldots,k$, consider the evaluation morphism $\theta
_{a}:\mathbb{K}[T]\rightarrow\mathbb{L}$ which associates to any
$P\in\mathbb{K}[T]$, the polynomial $P(h_{a})\in\Lambda_{(k)}\mathbb{=R}%
(h_{1},\ldots,h_{k})$. We know that $\{s_{\lambda}^{(k)}\mid\lambda
\in\mathcal{B}_{k}\}$ is a basis of $\Lambda_{(k)}$ over $\mathbb{R}$ thus
each power $h_{a}^{i},i\in\mathbb{Z}_{\geq0}$ decomposes on the basis of
$k$-Schur functions with real coefficients. By Corollary~\ref{reducedCase},
$h_{a}^{i}$ then decomposes on the family $\mathcal{I}$ with coefficients in
$\mathbb{K}$. Thus, $P(h_{a})$ also decomposes on the family $\mathcal{I}$
with coefficients in $\mathbb{K}$. Since $\mathcal{I}$ is a finite set, this
shows that $\Imrm (\theta_{a})$ is a finite-dimensional $\mathbb{K}%
$-subspace of $\mathbb{L}$, thus $h_{a}$ is algebraic over $\mathbb{K}$.

\eqref{prop4.2_2}:~Since $h_{1},\ldots,h_{k}$ are algebraic over $\mathbb{K}$, we get
$\mathbb{K}[h_{1},\ldots,h_{k}]=\mathbb{K}(h_{1},\ldots,h_{k})\subset
\mathbb{L}$. We also have $\mathbb{L}=\mathbb{R}(h_{1},\ldots,h_{k}%
)\subset\mathbb{K}(h_{1},\ldots,h_{k})$ since $\mathbb{K}$ is an extension of
$\mathbb{R}$. Thus, $\mathbb{L}=\mathbb{K}[h_{1},\ldots,h_{k}]$.

\eqref{prop4.2_3}:~This easily follows from \eqref{prop4.2_1} and \eqref{prop4.2_2}.

\eqref{prop4.2_4}:~By using the same arguments as in the proof of \eqref{prop4.2_1}, we get that each element
$Q(h_{1},\ldots,h_{k})$ with $Q\in\mathbb{K}[T_{1},\ldots,T_{k}]$ decomposes
on the family $\mathcal{I}$ with coefficients in $\mathbb{K}$. Assume we have
\[
\sum_{\kappa\mid k\text{-irreducible}}c_{\kappa}s_{\kappa}^{(k)}=0
\]
with $c_{\kappa}\in\mathbb{K}=\mathbb{R}(s_{R_{1}},\ldots,s_{R_{k}})$ for any
$k$-irreducible partition $\kappa$. Up to multiplication, we can assume these
coefficients belong in fact to $\mathbb{R}[s_{R_{1}},\ldots,s_{R_{k}}]$. Set
\[
c_{\kappa}=\sum_{\beta\in\mathbb{Z}_{\geq0}^{k}}a_{\beta}^{(\kappa)}s_{R_{1}%
}^{\beta_{1}}\cdots s_{R_{k}}^{\beta_{k}}%
\]
where all the coefficients $a_{\beta}^{(\kappa)}$ are equal to $0$ up to a
finite number (in which case $a_{\beta}^{(\kappa)}$ is real). We get%
\begin{equation}
\sum_{\kappa\mid k\text{-irreducible}}\sum_{\beta\in\mathbb{Z}_{\geq0}^{k}%
}a_{\beta}^{(\kappa)}s_{R_{1}}^{\beta_{1}}\cdots s_{R_{k}}^{\beta_{k}%
}s_{\kappa}^{(k)}=0. \label{CL}%
\end{equation}
By Remark~\ref{Rem_bij}, there is a bijection between the set of $k$-bounded
partitions and that of pairs $(\beta,\kappa)$ with $\kappa$ $k$-irreducible
and $\beta\in\mathbb{Z}_{\geq0}^{k}$. So equation~\eqref{CL} gives in fact a
linear combination of $k$-Schur functions over $\mathbb{R}$ equal to $0$.
Since we know that the set of $k$-Schur functions is a basis of $\Lambda
_{(k)}$, this imposes that each coefficient $a_{\beta}^{(\kappa)}$ is in fact
equal to $0$. So the family $\mathcal{I}$ is a $\mathbb{K}$-basis of
$\mathbb{L}$ and we have $[\mathbb{L}:\mathbb{K}]=\card (\mathcal{I}%
)=k!.$

\eqref{prop4.2_5}:~The characteristic polynomial of each $h_{a},a=1,\ldots,k$ belongs to
$\mathbb{A}[T]$ because the multiplication by $h_{a}$ on the basis
$\mathcal{I}$ makes only appear coefficients in $\mathbb{A}$. Thus, each
$h_{a}$ is an integral element of $\Lambda_{(k)}=\mathbb{A}[h_{1},\ldots
,h_{k}]$ over $\mathbb{A}.$
\end{proof}

\subsection{Primitive element}

\label{subsec_primitiveel}By Proposition~\ref{PropLK}, $s_{1}=h_{1}$ is
algebraic over $\mathbb{K}$. Denote by $\Pi$ its minimal polynomial. Observe
that $\Pi$ is an irreducible polynomial of $\mathbb{K}[T]$. Also write
$\boldsymbol{\Phi}$ for the matrix of the multiplication by $s_{1}$ in
$\mathbb{L}$ in the basis $\mathcal{I}=\{s_{\kappa}^{(k)}\mid\kappa
\in\mathcal{P}_{\irr }\}$ (here, we assume we have fixed once and for all a
total order on the set $\mathcal{P}_{\irr }$ of $k$-irreducible
partitions). Let $\Xi(T)=\det (TI_{k!}-\boldsymbol{\Phi})$ be the
characteristic polynomial of the matrix~$\boldsymbol{\Phi}$. We thus have that
$\Pi$ divides $\Xi$. Moreover both polynomials belong to $\mathbb{A}[T]$ since
the entries of the matrix $\boldsymbol{\Phi}$ are in the ring $\mathbb{A}$. We
have in fact the following stronger proposition.

\begin{prop}\label{Prop_InvSim}\ \\*[-1.2em]

\begin{enumerate}
\item\label{prop4.3_1} The invariant factors of the multiplication by $s_{1}$ are all equal to
$\Pi$: there exists an integer $m$ such that $\Xi=\Pi^{m}.$

\item\label{prop4.3_2} The coefficients of $\Pi$ and $\Xi$ are invariant under the flips of
$s_{R_{a}}$ and $s_{R_{k-a+1}}$ for any $a=1,\ldots,\left\lfloor \frac{k}%
{2}\right\rfloor $.
\end{enumerate}
\end{prop}

\begin{proof}
For~\eqref{prop4.3_1}, write $P_{1},\ldots,P_{m}$ for the invariant factors of $\boldsymbol{\Phi}%
$. We must have $P_{1}/P_{2}/\cdots/P_{m}$, $P_{m}=\Pi$ and $P_{1}P_{2}\cdots
P_{m}=\Xi$. Since $\Pi$ is irreducible, this imposes $P_{1}=\cdots=P_{m}=\Pi$
which gives the assertion.

For~\eqref{prop4.3_2}, we apply $\omega_{k}$ to each equality $\Pi(s_{(1)})=0$ and $\Xi
(s_{(1)})=0$.
\end{proof}

%\bigskip

\begin{exam}
\label{ex_k3_Phi}For $k=3,$ we get by listing the $k$-bounded partitions as
$\emptyset,(1),(2)$, $(1,1),(2,1)$ and $(2,1,1)$%
\[
\boldsymbol{\Phi}=\left(
\begin{matrix}
%[c]{cccccc}%
0 & 0 & s_{R_{1}} & s_{R_{3}} & s_{R_{2}} & 0\\
1 & 0 & 0 & 0 & 0 & s_{R_{2}}\\
0 & 1 & 0 & 0 & 0 & s_{R_{3}}\\
0 & 1 & 0 & 0 & 0 & s_{R_{1}}\\
0 & 0 & 1 & 1 & 0 & 0\\
0 & 0 & 0 & 0 & 1 & 0
\end{matrix}
\right)
\]
This gives $\Xi(T)=T^{6}-2\left( s_{R_{1}}+s_{R_{3}}\right) T^{3}-4s_{R_{2}%
}T^{2}+\left( s_{R_{1}}-s_{R_{3}}\right) ^{2}.$ Observe the symmetry of
$\boldsymbol{\Phi}$ which will be elucidated in \S~\href{https://alco.centre-mersenne.org/articles/10.5802/alco.147/}{6.3}.
\end{exam}

In Proposition~\ref{Prop_InvSim}, we have just used that $\boldsymbol{\Phi}$
is the matrix of the multiplication by $s_{(1)}$ which is algebraic over
$\mathbb{K}$. Given any morphism $\widetilde{\varphi}:\mathbb{A}%
\rightarrow\mathbb{R}$ such that $\widetilde{\varphi}(s_{R_{a}})\geq0$ for any
$a=1,\ldots,k$, the matrix $\widetilde{\varphi}(\boldsymbol{\Phi})$ obtained
by replacing in $\boldsymbol{\Phi}$ each rectangle Schur function $s_{R_{a}}$
by $\widetilde{\varphi}(s_{R_{a}})$ coincides with the matrix $\Phi$ defined
in \S~\ref{subsec_MatrixFI}.\ Thus $\Phi=\widetilde{\varphi}%
(\boldsymbol{\Phi})$ has nonnegative entries.\ We can now state the
main theorem of this section.

\begin{theorem}
We have $\mathbb{L}=\mathbb{K}(s_{(1)})$, that is $s_{(1)}$ is a primitive
element for $\mathbb{L}$ regarded as an extension of $\mathbb{K}$.
\end{theorem}

\begin{proof}
It suffices to show that $\Pi=\Xi$.\ By Proposition~\ref{Prop_InvSim}, we
already know that $\Xi=\Pi^{m}$ with $m\in\mathbb{Z}_{>0}$. Then by Frobenius
reduction, there exists an invertible matrix $P$ with entries in $\mathbb{K}$
such that
\begin{equation}
\boldsymbol{\Phi}=P\left(
\begin{matrix}
%[c]{ccc}%
\mathcal{C}_{\Pi} & \boldsymbol{0} & \boldsymbol{0}\\
\boldsymbol{0} & \ddots & \boldsymbol{0}\\
\boldsymbol{0} & \boldsymbol{0} & \mathcal{C}_{\Pi}%
\end{matrix}
\right) P^{-1}, \label{FIB}%
\end{equation}
that is, the matrix $\boldsymbol{\Phi}$ is equivalent to a block diagonal
matrix with $k$ blocks equal to $\mathcal{C}_{\Pi}$ the companion matrix of
the polynomial $\Pi$. By multiplying the columns of the matrix $P$ by elements
of $\mathbb{A}$, one can also assume that the entries of $P$ belong to
$\mathbb{A}$.\ Then we can write $P^{-1}=\frac{1}{\det(P)}Q$ where $Q$ has
also entries in $\mathbb{A}$ and $\det(P)\in\mathbb{A}$ is nonzero. Since
$\det(P)\in\mathbb{A}=\mathbb{R}[s_{R_{1}},\ldots,s_{R_{k}}]$ is nonzero,
there exists a nonzero polynomial $F\in\mathbb{R}[T_{1},\ldots,T_{k}]$ such
that $\det(P)=F(s_{R_{1}},\ldots,s_{R_{k}})$.\ Also a morphism $\widetilde
{\varphi}:\mathbb{A}\rightarrow\mathbb{R}$ such that $\widetilde{\varphi
}(s_{R_{a}})\geq0$ for any $a=1,\ldots,k$ is characterized by the datum of the
$\widetilde{\varphi}(s_{R_{a}})$'s. The polynomial $F$ being nonzero, one can
find $(r_{1},\ldots,r_{k})\in\mathbb{R}_{>0}^{k}$ such that $F(r_{1}%
,\ldots,r_{k})\neq0$.\ For such a $k$-tuple, let us define $\widetilde
{\varphi}$ by setting $\widetilde{\varphi}(s_{R_{a}})=r_{a}$. Then
$\widetilde{\varphi}(\det(P))\neq0$ and we can apply $\widetilde{\varphi}$ to~\eqref{FIB} which gives%
\[
\Phi=\widetilde{\varphi}(P)\left(
\begin{matrix}
%[c]{ccc}%
\mathcal{C}_{\widetilde{\varphi}(\Pi)} & \boldsymbol{0} & \boldsymbol{0}\\
\boldsymbol{0} & \ddots & \boldsymbol{0}\\
\boldsymbol{0} & \boldsymbol{0} & \mathcal{C}_{\widetilde{\varphi}(\Pi)}%
\end{matrix}
\right) \widetilde{\varphi}(P)^{-1}.
\]
The matrix $\Phi$ has nonnegative entries and is irreducible by Lemma
\ref{irredPhi}.\ So, by the Perron--Frobenius theorem, it admits a unique
eigenvalue $t>0$ of maximal modulus and the corresponding eigenspace is
one-dimensional. This eigenvalue $t$ should also be a root of $\widetilde
{\varphi}(\Pi)$, thus there is a vector $v\in\mathbb{R}^{d}$ with $d=\deg
(\Pi)$ such that $\mathcal{C}_{\widetilde{\varphi}(\Pi)}v=tv$. Then we get $m$
right eigenvectors of $\Phi$ linearly independent on $\mathbb{R}^{dm}$%
\[
\left(
\begin{matrix}
%[c]{c}%
\boldsymbol{v}\\
\boldsymbol{0}\\
\vdots\\
\boldsymbol{0}%
\end{matrix}
\right) ,\left(
\begin{matrix}
%[c]{c}%
\boldsymbol{0}\\
\boldsymbol{v}\\
\vdots\\
\boldsymbol{0}%
\end{matrix}
\right) ,\ldots,\left(
\begin{matrix}
%[c]{c}%
\boldsymbol{0}\\
\boldsymbol{0}\\
\vdots\\
\boldsymbol{v}%
\end{matrix}
\right) .
\]
Since the eigenspace considered is one-dimensional, this means that $m=1$ and
we are done.
\end{proof}

\begin{coro}
\label{Cor_Delta}There exist $\Delta\in\mathbb{A}$ and for each irreducible
$k$-bounded partition $\kappa$ a polynomial $P_{\kappa}\in\mathbb{A}[T]$ such
that
\[
s_{\kappa}^{(k)}=\frac{1}{\Delta}P_{\kappa}(s_{(1)}).
\]
In particular, for any morphism $\varphi:\Lambda_{(k)}\rightarrow\mathbb{R}$
such that $\varphi(\Delta)\neq0$ we have%
\[
\varphi(s_{\kappa}^{(k)})=\frac{1}{\varphi(\Delta)}\varphi(P_{\kappa}%
)(\varphi(s_{(1)})).
\]

\end{coro}

\begin{proof}
Since $s_{(1)}$ is a primitive element for $\mathbb{L}$ regarded as an
extension of $\mathbb{K}$, $\{1,s_{(1)},\ldots,s_{(1)}^{k!-1}\}$ is a
$\mathbb{K}$-basis of $\mathbb{L}$.\ It then suffices to consider the matrix
$M$ whose columns are the vectors $s_{(1)}^{i},i=0,\ldots,k!-1$ expressed on
the basis $\mathcal{I}=\{s_{\kappa}^{(k)},\kappa\in\mathcal{P}_{\irr %
}\}$. Its inverse can be written $M^{-1}=\frac{1}{\det M}N$ where the entries
of $N$ belong to $\mathbb{A}$. So we have $\Delta=\det(M)$ and the entries on
each columns of the matrix $N$ give the polynomials $P_{\kappa},\kappa
\in\mathcal{P}_{\irr }$.
\end{proof}

\begin{rema}
In fact we get the equality of $\mathbb{A}$-modules $\Lambda_{(k)}=\frac
{1}{\Delta}\mathbb{A}[s_{1}]$.\ In particular, the polynomial $\Delta$ (once
assumed monic) only depends on $\Lambda_{(k)}$ and $\mathbb{A}[s_{1}]$ and not
on the choice of the bases considered in these $\mathbb{A}$-modules. Indeed a
basis change will multiply $\Delta$ by an invertible element in $\mathbb{A}$,
that is by a nonzero real.
\end{rema}

\begin{exam}
For $k=2$ we get
\[
\boldsymbol{\Phi}=\left(
\begin{matrix}
%[c]{cc}%
0 & s_{R_{1}}+s_{R_{2}}\\
1 & 0
\end{matrix}
\right) \text{ and }M=I_{2}.
\]

\end{exam}

\begin{exam}
\label{Examk=3}For $k=3$ and with the same convention as in Example
\ref{ex_k3_Phi}, we get%
\begin{align*}
M&=\left(
\begin{matrix}
%[c]{cccccc}%
1 & 0 & 0 & s_{R_{1}}+s_{R_{3}} & 2s_{R_{2}} & 0\\
0 & 1 & 0 & 0 & s_{R_{1}}+s_{R_{3}} & 4s_{R_{2}}\\
0 & 0 & 1 & 0 & 0 & s_{R_{1}}+3s_{R_{3}}\\
0 & 0 & 1 & 0 & 0 & 3s_{R_{1}}+s_{R_{3}}\\
0 & 0 & 0 & 2 & 0 & 0\\
0 & 0 & 0 & 0 & 2 & 0
\end{matrix}
\right) \\
\intertext{and}
M^{-1}&=\left(
\begin{matrix}
%[c]{cccccc}%
1 & 0 & 0 & 0 & -\frac{s_{R_{1}}+s_{R_{3}}}{2} & -s_{R_{2}}\\
0 & 1 & \frac{2s_{R_{2}}}{s_{R_{1}}-s_{R_{3}}} & \frac{2s_{R_{2}}}{s_{R_{3}%
}-s_{R_{1}}} & 0 & -\frac{s_{R_{1}}+s_{R_{3}}}{2}\\
0 & 0 & \frac{3s_{R_{1}}+s_{R_{3}}}{2s_{R_{1}}-2s_{R_{3}}} & \frac{s_{R_{1}%
}+3s_{R_{3}}}{2s_{R_{3}}-2s_{R_{1}}} & 0 & 0\\
0 & 0 & 0 & 0 & \frac{1}{2} & 0\\
0 & 0 & 0 & 0 & 0 & \frac{1}{2}\\
0 & 0 & \frac{-1}{2s_{R_{1}}-2s_{R_{3}}} & \frac{-1}{2s_{R_{3}}-2s_{R_{1}}} &
0 & 0
\end{matrix}
\right) %\allowbreak.
\end{align*}
So in particular, $s_{(2,1,1)}^{(3)}=\frac{1}{2}s_{(1)}^{4}-\frac{1}%
{2}(s_{R_{1}}+s_{R_{3}})s_{(1)}-s_{R_{2}}$.
\end{exam}

\subsection{Algebraic variety associated to fixed values of rectangles}

Recall that $\Lambda_{(k)}=\mathbb{R}[h_{1},\ldots,h_{k}]$ and each rectangle
Schur polynomial can be written $s_{R_{a}}=JT_{a}(h_{1},\ldots,h_{k})$ for any
$a=1,\ldots,k$ where $JT_{a}\in\mathbb{R}[h_{1},\ldots,h_{k}]$ is given by the
Jacobi--Trudi determinantal formula. Consider $\vec{r}=(r_{1},\ldots,r_{k}%
)\in\mathbb{R}_{\geq0}^{k}$.

\begin{defi}
\label{Def_Variety}Let $\mathcal{R}_{\vec{r}}$ be the algebraic variety of
$\mathbb{R}^{k}$ defined by the equations $s_{R_{a}}=r_{a}$ for any
$a=1,\ldots,k$.
\end{defi}

We can consider the algebra $\overline{\Lambda}_{(k)}:=\Lambda_{(k)}/J$ where
$J$ is the ideal generated by the relations $s_{R_{a}}=r_{a}$ for any
$a=1,\ldots,k$. Write $\overline{\varphi}:\Lambda_{(k)}\rightarrow
\Lambda_{(k)}/J$ for the canonical projection obtained by specializing in
$\Lambda_{(k)}$ each rectangle Schur function $s_{R_{a}}$ to $r_{a}$.\ We
shall write for short $\overline{b}=\overline{\varphi}(b)$ for any
$b\in\Lambda_{(k)}$. Clearly $\overline{\Lambda}_{(k)}=\Lambda_{(k)}/J$ is a
finite-dimensional $\mathbb{R}$-algebra and $\overline{\Lambda}_{(k)}%
=\vect \langle\overline{s}_{\kappa}\mid\kappa$ irreducible$\rangle$.
The following proposition shows that the non-cancellation of $\Delta$ can be
naturally interpreted as a condition for the multiplication by $\overline
{s}_{1}$ to be a cyclic morphism in $\overline{\Lambda}_{(k)}$.

\begin{prop}
\label{primitive_Specialization} \ 

\begin{enumerate}
\item\label{prop4.11_1} The algebra $\overline{\Lambda}_{(k)}$ has dimension $k!$ over
$\mathbb{R}$ and $\{\overline{s}_{\kappa}\mid\kappa\in\mathcal{P}%
_{\irr }\}$ is a basis of $\overline{\Lambda}_{(k)}$.

\item\label{prop4.11_2} We have $\overline{\Lambda}_{(k)}=\mathbb{R}[\overline{s}_{1}]$ if and
only if $\overline{\Delta}\neq0$.
\end{enumerate}
\end{prop}

\begin{proof}
\eqref{prop4.11_1}:~For any $k$-bounded partition $\lambda$, write $\lambda=R_{1}^{m_{1}}%
\sqcup\cdots\sqcup R_{k}^{m_{k}}\sqcup\kappa(\lambda)$ for its decomposition into
rectangles and irreducible partition, and set $u(\lambda)=r_{1}^{m_{1}}\cdots
r_{k}^{m_{k}}$.\ If all $m_{i}$ are equal to zero for $1\leq i\leq k$, we set $u(\lambda)=0$. Let us prove that $J$ regarded as a $\mathbb{R}$-vector space has basis
$\{s_{\lambda}^{(k)}-u(\lambda)s_{\kappa(\lambda)}^{(k)}\mid\lambda\in \mathcal{B}_{k}\setminus \mathcal{P}_{\irr }\}$. First, the latter set is linearly independent and is thus a basis of a vector subspace $V\subset \Lambda_{(k)}$. It remains to prove that $V=J$. 

Let $\lambda=R_{1}^{m_{1}}%
\sqcup\cdots\sqcup R_{k}^{m_{k}}\sqcup\kappa(\lambda)$ be a reducible $k$-bounded partition and let $\kappa'$ be a $k$-bounded partition. Then,
\begin{align*}
s_{\kappa'}^{(k)} (s_{\lambda}^{(k)}-u(\lambda)s_{\kappa(\lambda)}^{(k)})&=s_{\kappa'}^{(k)}s_{\kappa(\lambda)}^{(k)}\left(s_{R_{1}^{m_{1}}%
\sqcup\cdots\sqcup R_{k}^{m_{k}}}^{(k)}-u(R_{1}^{m_{1}}%
\sqcup\cdots\sqcup R_{k}^{m_{k}})\right)\\
&=\sum_{\nu\in \mathcal{B}_{k}}c_{\kappa(\lambda),\kappa'}^{\nu(k)} s_{\nu}^{(k)}\left(s_{R_{1}^{m_{1}}%
\sqcup\cdots\sqcup R_{k}^{m_{k}}}^{(k)}-u(R_{1}^{m_{1}}%
\sqcup\cdots\sqcup R_{k}^{m_{k}})\right)
\end{align*}
Let us write $\lambda'=R_{1}^{m_{1}}%
\sqcup\cdots\sqcup R_{k}^{m_{k}}$. For each $\nu\in\mathcal{B}_{k}$ written $\nu=R_{1}^{n_{1}}%
\sqcup\cdots\sqcup R_{k}^{n_{k}}\sqcup\kappa(\nu):=\nu'\sqcup \kappa(\nu)$, we have 
\begin{align*}
s_{\nu}^{(k)}(s_{\lambda'}^{(k)}-u(\lambda'))&=s_{\kappa(\nu)}^{(k)}(s_{\lambda'\sqcup \nu'}^{(k)}-u(\lambda')s_{\nu'}^{(k)})\\
&=s_{\kappa(\nu)}^{(k)}\left(s_{\lambda'\sqcup \nu'}^{(k)}-u(\lambda'\sqcup \nu')+u(\lambda'\sqcup \nu')-u(\lambda')s_{\nu'}^{(k)}\right)\\
&=s_{\kappa(\nu)}^{(k)}\left(s_{\lambda'\sqcup \nu'}^{(k)}-u(\lambda'\sqcup \nu')-u(\lambda')(s_{\nu'}^{(k)} -u(\nu'))\right)\\
&=s_{\lambda'\sqcup \nu'\sqcup \kappa(\nu)}^{(k)}-u(\lambda'\sqcup \nu'\sqcup \kappa(\nu))s_{\kappa(\nu)}^{(k)}-u(\lambda')(s_{\nu}^{(k)}-u(\nu)s_{\kappa(\nu)}^{(k)})\in V,
\end{align*}
where we have used $u(\lambda'\sqcup \nu')=u(\lambda'\sqcup \nu'\sqcup \kappa(\nu))$ on the last equality. Hence, $s_{\kappa'}^{(k)} (s_{\lambda}^{(k)}-u(\lambda)s_{\kappa(\lambda)}^{(k)})\in V$ and $V$ is an ideal of $\Lambda_{(k)}$. Since all generators $s_{R_{a}}-r_{a}$ of $J$ are in $V$, we have in particular $J\subset V$. 

Let us show that $(s_{\lambda}^{(k)}-u(\lambda)s_{\kappa(\lambda)}^{(k)})\in J$ for all reducible $k$-bounded partitions $\lambda=R_{1}^{m_{1}}%
\sqcup\cdots\sqcup R_{k}^{m_{k}}\sqcup\kappa(\lambda)$. Since $(s_{\lambda}^{(k)}-u(\lambda)s_{\kappa(\lambda)}^{(k)})=s_{\kappa(\lambda)}^{(k)}(s_{R_{1}^{m_{1}}%
\sqcup\cdots\sqcup R_{k}^{m_{k}}}^{(k)}-u(R_{1}^{m_{1}}%
\sqcup\cdots\sqcup R_{k}^{m_{k}}))$, we just have to show that $s_{R_{1}^{m_{1}}%
\sqcup\cdots\sqcup R_{k}^{m_{k}}}^{(k)}-u(R_{1}^{m_{1}}%
\sqcup\cdots\sqcup R_{k}^{m_{k}})\in J$ for all tuple $\vec{m}=(m_{1},\ldots,m_{k})$ different from zero. We prove this by induction on $l(\vec{m})=\sum_{i}^k m_{i}$. For $\vec{m}$ such that $l(\vec{m})=1$, this is true. Suppose that $\vec{m}$ is such that $l(\vec{m})>1$ and assume without loss of generality that $m_{1}\geq 1$. Then,
\begin{multline*}
s_{R_{1}^{m_{1}}%
\sqcup\cdots\sqcup R_{k}^{m_{k}}}^{(k)} -u(R_{1}^{m_{1}}%
\sqcup\cdots\sqcup R_{k}^{m_{k}})=s_{R_{1}} (s_{R_{1}^{m_{1}-1}%
\sqcup\cdots\sqcup R_{k}^{m_{k}}}^{(k)} -u(R_{1}^{m_{1}-1}%
\sqcup\cdots\sqcup R_{k}^{m_{k}}))\\
+u(R_{1}^{m_{1}-1}%
\sqcup\cdots\sqcup R_{k}^{m_{k}})(s_{R_{1}}-u(R_{1}))\in J,
\end{multline*}
where we have used the induction hypothesis to prove that the first term of the right hand side is in $J$. Hence, $V\subset J$, and finally $V=J$.

Observe also that $\{s_{\lambda}^{(k)}-u(\lambda)s_{\kappa(\lambda)}^{(k)}\mid\lambda$ $k$-bounded partition$\}$ is a basis of $\Lambda_{(k)}$, since it is a triangular change of basis from the standard basis $\{ s_{\lambda}^{(k)}\mid \lambda$ $k$-bounded partition$\}$. Hence, we have a decomposition 
\[
\Lambda_{(k)}=\vect \langle s_{\lambda}^{(k)}-u(\lambda)s_{\kappa(\lambda)}^{(k)}\mid\lambda\in \mathcal{P}_{\irr }\rangle\oplus J.
\]
We deduced that a basis of $\Lambda_{(k)}/J$ is given by $\{ \overline{\varphi}(s_{\lambda}^{(k)}-u(\lambda)s_{\kappa(\lambda)}^{(k)})\mid\lambda\in\mathcal{P}_{\irr }\}$. Since $u(\lambda)=0$ and $\lambda=\kappa(\lambda)$ when $\lambda$ is irreducible, we get that $\{ \overline{s}_{\kappa}\mid\kappa\in\mathcal{P}_{\irr }\}$ is a basis of $\Lambda_{(k)}/J$.

\eqref{prop4.11_2}:~We have $\overline{\Delta}\neq0$ if and only if $\{\overline{s}_{1}^{r}%
\mid0\leq r\leq k!-1\}$ is a basis of $\overline{\Lambda}_{(k)}$ since
$\{\overline{s}_{\kappa}\mid\mathcal{P}_{\irr }\}$ is a basis of
$\overline{\Lambda}_{(k)}$ and $\overline{\Delta}$ is then the determinant
between the two bases.
\end{proof}

The following proposition is classical, we prove it for completion.

\begin{prop}
\label{Prop_Variety}The algebraic variety $\mathcal{R}_{\vec{r}}$ is finite.
\end{prop}

\begin{proof}
It suffices to see that the algebraic variety $\mathcal{R}_{\vec{r}%
}^{\mathbb{C}}$ of $\mathbb{C}^{k}$ defined by the equations $s_{R_{a}}=r_{a}$
for any $a=1,\ldots,k$ is finite.\ We can decompose $\mathcal{R}_{\vec{r}%
}^{\mathbb{C}}=V_{1}\cup\cdots\cup V_{m}$ into its irreducible components.\ To
each such component $V_{j}$ is associated a prime ideal $J_{j}$ and we have
$J=J_{1}\cap\cdots\cap J_{m}$. Therefore for any $j=1,\ldots,m,$
$\mathbb{C}[h_{1},\ldots,h_{k}]/J_{j}$ is a finite-dimensional algebra
(because $\mathbb{C}[h_{1},\ldots,h_{k}]/J$ is) which is an integral domain.
So $\mathbb{C}[h_{1},\ldots,h_{k}]/J_{j}$ is in fact a field and $J_{j}$ is
maximal in $\mathbb{C}[h_{1},\ldots,h_{k}]$. By using Hilbert's
Nullstellensatz, we obtain that each $V_{j}$ reduces to a point, so
$\mathcal{R}_{\vec{r}}^{\mathbb{C}}$ is finite.
\end{proof}

\section{Nonnegative morphisms on \texorpdfstring{$\Lambda_{(k)}$}{Lambda k}}\label{section5}

\subsection{Nonnegative morphisms with \texorpdfstring{$\Phi$}{Phi} irreducible}

We show now that when the matrix $\Phi$ introduced in
\S~\ref{subsec_MatrixFI} is irreducible, the values of $\varphi$ on the
rectangle Schur functions $s_{R_{a}},$ $1\leq a\leq k$ determine completely
the morphism $\varphi$. Recall we have denoted by $\mathbb{A}=\mathbb{R}%
[s_{R_{1}},\ldots,s_{R_{k}}]$ the subalgebra of $\Lambda_{(k)}=\mathbb{R}%
[h_{1},\ldots,h_{k}]$ generated by the $k$-rectangle Schur functions. Also we
have $\mathcal{I}=\{s_{\kappa}^{(k)}\mid\kappa\in\mathcal{P}_{\irr }%
\}$. Set $\mathcal{R}=\{s_{R_{a}}\}_{1\leq a\leq k}$.

\begin{theorem}
\label{principalDomain} \ 

\begin{enumerate}
\item\label{theo5.1_1} Let $\varphi:\mathbb{A}\rightarrow\mathbb{R}$ be a morphism, nonnegative
on $\mathcal{R}$ and such that its associated matrix $\Phi$ is irreducible.
Then there exists a unique morphism $\tilde{\varphi}:\Lambda_{(k)}%
\longrightarrow\mathbb{R}_{\geq0}$ extending $\varphi$, nonnegative on the
$k$-Schur functions and positive on $\mathcal{I}$.

\item\label{theo5.1_2} A positive morphism $\varphi:\Lambda_{(k)}\longrightarrow\mathbb{R}$ is
uniquely determined by its values on~$\mathcal{R}$.
\end{enumerate}
\end{theorem}

\begin{proof}
\eqref{theo5.1_1}: Let us prove the existence of $\tilde{\varphi}$. Set $(r_{1},\ldots,
r_{k})=(\varphi(s_{R_{1}}),\ldots,\varphi(s_{R_{k}}))$, and assume first that
$\Delta(r_{1},\ldots,r_{k})\not =0$ where $\Delta$ is the polynomial defined
in Corollary~\ref{Cor_Delta}. We have to show that there exists a morphism
$\widetilde{\varphi}$ on $\Lambda_{(k)}$ such that $\widetilde{\varphi}$ is
positive on $\mathcal{I}$ and $\widetilde{\varphi}(s_{R_{a}})=r_{a}$ for
$a=1,\ldots,k$. The set of morphisms from $\Lambda_{(k)}$ to $\mathbb{R}$
which takes values $r_{a}$ on $s_{R_{a}}$ for $a=1\ldots k$ is in bijection
with the set of morphisms from $\overline{\Lambda}_{(k)}$ to $\mathbb{R}.$
Here recall that $\overline{\Lambda}_{(k)}=\Lambda_{(k)}/J$ with $J$ the ideal
generated by the relations $s_{R_{a}}=r_{a}$ for $a=1,\ldots,k$. Since we have
assumed $\bar{\Delta}\not =0$, Proposition~\ref{primitive_Specialization}
yields that $\overline{\Lambda}_{(k)}=\mathbb{R}[\overline{s}_{1}]$. There
exists one morphism from $\mathbb{R}[\overline{s}_{1}]$ to $\mathbb{R}$ for
each real root of the minimal polynomial of $\overline{s}_{1}$, which is
$\overline{\Xi}$ because $\deg(\overline{\Xi})=k!=\dim(\overline{\Lambda
}_{(k)})$. Let $t$ be the root of $\overline{\Xi}$ with maximal modulus. It is
positive since $\overline{\Xi}$ is the characteristic polynomial of the
irreducible matrix $\Phi_{(r_{1},\ldots,r_{k})}$, and moreover it is the
Perron--Frobenius eigenvalue of $\Phi$. Then, the specialization $\overline
{s}_{1}=t$ yields a morphism from $\overline{\Lambda}_{(k)}$ to $\mathbb{R}$,
and by extension a morphism $\widetilde{\varphi}$ from $\Lambda_{(k)}$ to
$\mathbb{R}$. In particular, the equality~\eqref{PieriPhi} holds for
$\widetilde{\varphi}$, and the vector $X=(\widetilde{\varphi}(s_{\lambda
}^{(k)}))_{\lambda\in\mathcal{P}_{\irr }}$ is the eigenvector of $\Phi$
corresponding to the Perron--Frobenius eigenvalue $t$ and such that
$X(\emptyset)=1$. Therefore, $\widetilde{\varphi}$ is positive on
$\mathcal{I}$. So we have proved the existence of an extension of $\varphi$
with the right properties in the case where $\Delta\neq0$.

We now drop the hypothesis $\Delta=0$. Consider $\vec{r}=(r_{1},\ldots,r_{k})$
such that the matrix $\Phi$ is irreducible. For any extension $\widetilde
{\varphi}$ (if any) of $\varphi$ positive on $\mathcal{I}$, the 
equality~\eqref{PieriPhi} still holds. Therefore, $X^{\vec{r}}=(\widetilde{\varphi
}(s_{\lambda}^{(k)})_{\lambda\in\mathcal{P}_{\irr }})$ should then be
the eigenvector of $\Phi$ corresponding to the Perron--Frobenius eigenvalue
$t_{\vec{r}}$ and such that $X^{\vec{r}}(\emptyset)=1$. For $\mu\in
\mathcal{B}_{k}$, set
\begin{equation}
\label{expression_morphism}\widetilde{\varphi}_{\vec{r}}(s_{\mu}^{(k)}%
)=\prod_{a=1}^{k}r_{a}^{p_{a}^{\mu}}X^{\vec{r}}(\widetilde{\mu})\text{ with
}s_{\mu}^{(k)}=\prod_{a=1}^{k}s_{R_{a}}^{p_{a}^{\mu}}s_{\widetilde{\mu}}%
^{(k)}\text{ , }\tilde{\mu}\text{ irreducible.}%
\end{equation}
Then, the map $\widetilde{\varphi}_{\vec{r}}$ is an extension of $\varphi$
which is positive on $\mathcal{I}$ and nonnegative on the $k$-Schur functions
by construction.\ So it just remains to prove that $\widetilde{\varphi}%
_{\vec{r}}$ is a morphism.

Since $t_{\vec{r}}$ and $X^{\vec{r}}$ are continuous functions of $\vec{r}$ on
the set of irreducible matrices, the map $\widetilde{\varphi}_{\vec{r}}$ is a
continuous function of $\vec{r}$. The hypersurface $V(\Delta):=\{\Delta
(\vec{r})=0\}$ is closed in the Zariski topology, thus $V(\Delta)$ has empty
interior in the set $\Theta=\{\vec{r}\in\mathbb{R}_{\geq0}^{k}\mid\Phi
_{(r_{1},\ldots,r_{k})}$ is irreducible$\}$. By the previous arguments, for
all $\vec{r}\in\Theta$ outside $V(\Delta)$, the map $\widetilde{\varphi}%
_{\vec{r}}$ is a morphism and $\vec{r}\mapsto\widetilde{\varphi}_{\vec{r}}$ is
continuous on $\Theta$, thus $\widetilde{\varphi}_{\vec{r}}$ is a morphism for
$\vec{r}\in\overline{\Theta\setminus V(\Delta)}$. By Proposition
\ref{irredPhi}, the set $\Theta$ is open.\ Let $\vec{r}\in\Theta\cap
V(\Delta)$.\ Since the interior of $V(\Delta)$ is empty, one can define a
sequence $\vec{r}^{(n)}\in\Theta\setminus V(\Delta)$ which tends to $\vec
{r}\in\overline{\Theta\setminus V(\Delta)}$ as $n$ goes to infinity. Finally
we get that $\widetilde{\varphi}_{\vec{r}}$ is a morphism, as desired.

We now prove the uniqueness of the extension $\widetilde{\varphi}$. Let
$\widehat{\varphi}$ be another real extension of $\varphi$ positive on
$\mathcal{I}$. By~\eqref{PieriPhi}, the vector $\widehat{X}=(\widehat{\varphi
}(s_{\lambda}^{(k)}))_{\lambda\in\mathcal{P}_{\irr }}$ is also a left
eigenvector of $\Phi$ for the positive eigenvalue $\widehat{\varphi}(s_{1})$
normalized so that $\widehat{X}_{\emptyset}$ is equal to one. Since
$\widehat{X}$ has positive entries, it is the Perron--Frobenius eigenvector of
$\Phi$ and is thus equal to the vector $X^{\vec{r}}$ defined above. By the
morphism property of $\widehat{\varphi}$, for $\mu\in\mathcal{B}_{k}$ such
that $s_{\mu}^{(k)}=\prod_{a=1}^{k}s_{R_{a}}^{p_{a}^{\mu}}s_{\widetilde{\mu}%
}^{(k)}$ with $\tilde{\mu}$ irreducible, we have
\begin{equation}
\widehat{\varphi}(s_{\mu}^{(k)})=\prod_{a=1}^{k}\varphi(s_{R_{a}}^{p_{a}^{\mu
}})\widehat{X}(\widetilde{\mu})=\prod_{a=1}^{k}r_{a}^{p_{a}^{\mu}}%
X(\widetilde{\mu})=\widetilde{\varphi}(s_{\mu}^{(k)}),
\end{equation}
where we have used the definition of $\widetilde{\varphi}$ given in~\eqref{expression_morphism}. 
Therefore, $\widetilde{\varphi}=\widehat{\varphi
}$, and we have proven the uniqueness of the extension of $\varphi$ satisfying
the properties of the statement.

\eqref{theo5.1_2}: Let $\varphi$ be a positive morphism on $\Lambda_{(k)}$. Then, the
associated matrix $\Phi$ is irreducible by Proposition~\ref{irredPhi}. Hence,
the first part of the proposition yields that $\varphi$ is uniquely determined
by its values on $\mathcal{R}$.
\end{proof}

An immediate consequence of the latter theorem is the description of positive
extremal harmonic functions. We define an action of $\mathbb{R}_{>0}$ on
$\mathcal{F}(\mathcal{B}_{k},\mathbb{R}_{\geq0})$ by
\[
t\cdot\varphi(s_{\lambda}^{(k)})=t^{\left\vert \lambda\right\vert }%
\varphi(s_{\lambda}^{(k)}),
\]
for $t>0$, $\varphi\in\mathcal{F}(\mathcal{B}_{k},\mathbb{R}_{\geq0})$.

\begin{coro}\ \\*[-1.2em] 

\begin{enumerate}
\item Let $\varphi\in\partial\mathcal{H}^{+}(\mathcal{B}_{k})$, and suppose
that $\varphi$ is positive on the $k$-Schur functions. Then, $\varphi$ is
uniquely determined by its values on the $s_{R_{a}}$, $1\leq a\leq k$.

\item Assume the matrix $\Phi$ associated to $\varphi$ is irreducible. Then
there exists a unique $t>0$ such that $t^{-1}\cdot\varphi$ can be extended to
an element of $\partial\mathcal{H}^{+}(\mathcal{B}_{k})$ positive on~$\mathcal{I}$.
\end{enumerate}
\end{coro}

\begin{proof}
The only non immediate statement is the second one. Suppose that $\Phi$ is
irreducible. Then, by Theorem~\ref{principalDomain}, $\varphi$ can be extended
in a unique way to a morphism $\widetilde{\varphi}$ nonnegative on
$\mathcal{B}_{k}$ and positive on $\mathcal{I}$. Let $t=\widetilde{\varphi
}(s_{1})>0$. Then, $t^{-1}\cdot\widetilde{\varphi}$ is a nonnegative morphism
on $\mathcal{B}_{k}$ such that $t^{-1}\cdot\widetilde{\varphi}(s_{1})=1$,
which belongs to $\partial\mathcal{H}^{+}(\mathcal{B}_{k})$. It is clear that
$t^{-1}\cdot\widetilde{\varphi}$ extends $t^{-1}\cdot\varphi$. Also if
$\widehat{\theta}\in\partial\mathcal{H}^{+}(\mathcal{B}_{k})$ extends
$s^{-1}\cdot\varphi$ with $s>0$ and is positive on $\mathcal{I}$, then the
vector with entries $s\cdot\widehat{\theta}(s_{\lambda}^{(k)}),\lambda
\in\mathcal{P}_{\irr }$ is an eigenvector for $\Phi$ associated to the
eigenvalue $s$. Since it has positive entries, we get $t=s$ and $\widehat
{\theta}=t^{-1}\cdot\widetilde{\varphi}$.
\end{proof}

\begin{rema}
\label{Rem_param}It follows also from Proposition~\ref{irredPhi} and Assertion~\eqref{theo5.1_2}
 of Theorem~\ref{principalDomain}, that each morphism $\varphi:\mathbb{A}%
\rightarrow\mathbb{R}$ positive on $\mathcal{R}$ can be extended in a unique
way to a morphism $\widetilde{\varphi}:\Lambda_{(k)}\longrightarrow
\mathbb{R}_{\geq0}$ positive on the $k$-Schur functions. Also clearly, two
distinct such morphisms on $\mathbb{A}$ will yield distinct extensions on
$\Lambda_{(k)}$. Thus the morphisms $\widetilde{\varphi}:\Lambda
_{(k)}\longrightarrow\mathbb{R}_{\geq0}$ positive on the $k$-Schur functions
are parametrized by $\mathbb{R}_{>0}^{k}$ via the map which associates to each
such morphism its values on $\mathcal{R}$.
\end{rema}

\subsection{Two parametrizations of the positive morphisms}

\label{SubSec_Twoparam}The more immediate parametrization of the positive
morphisms $\varphi:\Lambda_{(k)}\rightarrow\mathbb{R}$ such that
$\varphi(s_{\lambda}^{(k)})>0$ for any $k$-bounded partition is obtained from
the factorization property (Corollary~\ref{reducedCase}) of the $k$-Schur
functions. Consider
\[
V=\left\{ (h_{1},\ldots,h_{k})\in\mathbb{R}^{k}\left\vert
\begin{array}[c]{ll}
%
JT_{R_{i}}(h_{1},\ldots,h_{k})>0,&i=1,\ldots,k\\[3pt]
JT_{\kappa}(h_{1},\ldots,h_{k})>0, &\forall\kappa\in\mathcal{P}_{\irr }%
\end{array}\right.
\right\} \subset\mathbb{R}_{>0}^{k}%
\]
where we have set $s_{Ri}=JT_{R_{i}}(h_{1},\ldots,h_{k})$ and $s_{\kappa
}=JT_{\kappa}(h_{1},\ldots,h_{k})$ where $JT_{R_{1}},\ldots,JT_{R_{k}}$ and
$JT_{\kappa},\kappa\in\mathcal{P}_{\irr }$ are polynomials in
$\mathbb{R}[X_{1},\ldots,X_{k}]$. To each point in $V$ corresponds a unique
positive morphism $\varphi$ defined on $\Lambda_{(k)}$. Now define%
\[
U=\{\vec{r}=(r_{1},\ldots,r_{k})\in\mathbb{R}_{>0}^{k}\}.
\]
As explained in Remark~\ref{Rem_param}, the positive morphisms $\varphi
:\Lambda_{(k)}\rightarrow\mathbb{R}$ such that $\varphi(s_{\lambda}^{(k)})>0$
for any $\lambda\in\mathcal{B}_{k}$ are parametrized by $U$ and we can define
a map $f:U\rightarrow V$ such that
\[
f(r_{1},\ldots,r_{k})=(\varphi(h_{1}),\ldots,\varphi(h_{k})).
\]
The map $f$ is then continuous on $U$ since the entries of the matrix $\Phi$
are and so is its Perron--Frobenius vector normalized at $1$ on $s_{\emptyset}%
$. Moreover, the map $f$ is bijective by Theorem~\ref{principalDomain} and we
have
\[
f^{-1}:\left\{
\begin{aligned}
%[c]{c}%
V&\rightarrow U\\
(h_{1},\ldots,h_{k})&\mapsto(JT_{R_{1}}(h_{1},\ldots,h_{k}),\ldots,JT_{R_{k}%
}(h_{1},\ldots,h_{k}))
\end{aligned}
\right.
\]
where the polynomials $JT_{R_{1}},\ldots,JT_{R_{k}}$ are given by the
Jacobi--Trudi determinantal formulas. In particular $f^{-1}$ is continuous on
$V$.

\begin{lemma}
\label{Lemma_Bounded}The map $f$ is bounded on any bounded subset of $U$.
\end{lemma}

\begin{proof}
Let $B\subset U$ be a bounded subset of $U$. By definition of $f$, for any
$\vec{r}=(r_{1},\ldots,r_{k})$ in $B$, $\varphi(h_{1})$ is the first
coordinate of $f(r_{1},\ldots,r_{k})$ and coincides with the Perron--Frobenius
eigenvalue of the matrix $\Phi$, that is with its spectral radius. Since the
spectral radius of a real matrix is a bounded function of its entries, we get
that $\varphi(h_{1})$ is bounded when $\vec{r}$ runs over $B$. To conclude,
observe that for any $a=2,\ldots,k$ we have $\varphi(h_{a})\leq\varphi
(h_{1})^{a}$ because $h_{1}=s_{1},$ the map $\varphi$ is multiplicative and
$h_{a}$ appears in the decomposition of $s_{1}^{a}$ on the basis of $k$-Schur
functions (which only makes appear nonnegative real coefficients).
\end{proof}

Now set
\begin{align*}
\overline{U} & =\{\vec{r}=(r_{1},\ldots,r_{k})\in\mathbb{R}_{\geq0}%
^{k}\}\\
\intertext{and }
\overline{V} & =\left\{ (h_{1},\ldots,h_{k})\in\mathbb{R}^{k}%\mid
\left\vert
\begin{array}[c]{ll}
JT_{R_{i}}(h_{1},\ldots,h_{k})\geq0,&i=1,\ldots,k\\[3pt]
JT_{\kappa}(h_{1},\ldots,h_{k})\geq0, &\forall\kappa\text{ }%
k\text{-irreducible}%
\end{array}
\right. 
\right\} \subset\mathbb{R}_{\geq0}^{k}.%
\end{align*}
Since $JT_{R_{1}},\ldots,JT_{R_{k}}$ are polynomials, we can extend $f^{-1}$
by continuity on $\overline{V}$ and get a continuous map $g:\overline
{V}\rightarrow\overline{U}$. But this is not immediate right now that $g$ is
bijective and $f$ can also be extended to a bijective map from $\overline{U}$
to $\overline{V}$. Observe nevertheless that if we can extend $f$ by
continuity on $\overline{U}$, the continuity of $g$ and $f$ will imply that
$f\circ g=id_{\overline{V}}$ and $g\circ f=id_{\overline{U}}.$ Therefore, to
extend $f$ by continuity it will suffice to prove that $\overline{U} $ and
$\overline{V}$ are homeomorphic by $f$.

\subsection{Extension of the map on \texorpdfstring{$\overline{U}$}{U}}

Let $\vec{r}_{0}\in\mathbb{R}_{\geq0}^{k}$, and denote by $A(\vec{r}_{0})$ the
set of limiting values of $f(\vec{r})$ as $\vec{r}$ goes to $\vec{r}_{0}$ in $\overline{U}=\mathbb{R}^{k}$.
Recall the notation of the previous paragraph, in particular the function $g$
is defined and continuous on $\overline{V}$ and $g=f^{-1}$ on $f(U)$.

\begin{lemma}
The set $A(\vec{r}_{0})$ is a connected subset of $\mathcal{R}_{\vec{r}_{0}}$
(see Definition~\ref{Def_Variety}).
\end{lemma}

\begin{proof}
Consider $K_{n}=B\left( \vec{r}_{0},\frac{1}{n}\right) \cap\mathbb{R}%
_{>0}^{k}$. This is a system of decreasing bounded connected neighborhoods of
$\vec{r}_{0}$ in $\mathbb{R}_{>0}^{k}$. By definition, $A(\vec{r}_{0}%
)=\bigcap_{n\geq1}\overline{f(K_{n})}$. By Lemma~\ref{Lemma_Bounded}, we know
that $f$ is bounded on bounded subsets of $U=\mathbb{R}_{>0}^{k}$, therefore
we get that $f(K_{n})$ is bounded and thus $\overline{f(K_{n})}$ is compact.
Since $f$ is continuous on $U$ and $K_{n}$ is connected, $f(K_{n})$ is also
connected, which implies that $\overline{f(K_{n})}$ is connected. Hence,
$A(\vec{r}_{0})$ is a decreasing intersection of connected compact sets, and
thus $A(\vec{r}_{0})$ is connected.

Let $\vec{h}\in A(\vec{r}_{0})$. Then, there exists a sequence $(\vec
{r_{n}})_{n\geq1}$ in $U$ converging to $\vec{r}_{0}$ such that $\vec{h}%
_{n}:=f(\vec{r_{n}})$ converges to $\vec{h}$ as $n$ goes to infinity. Since $g=f^{-1}$ on $\mathbb{R}_{>0}^{k}$, $g(\vec{h}_{n})=g\circ f(\vec
{r_{n}})=\vec{r_{n}}$ for $n\geq1$. Moreover, since $g$ is continuous and
$(\vec{h}_{n})_{n\geq1}$ converges to $\vec{h}$ as $n$ goes to infinity,
\[
g(\vec{h})=\lim_{n\rightarrow\infty}g(\vec{h}_{n})=\lim_{n\rightarrow\infty
}\vec{r_{n}}=\vec{r}_{0}%
\]
which implies that $\vec{h}\in\mathcal{R}_{\vec{r}_{0}}$.
\end{proof}


\begin{theorem}\label{Th_UVhomeo}\ \\*[-1.2em]

\begin{enumerate}
\item\label{theo5.6_1} The map $f$ is an homeomorphism from $\overline{U}$ to $\overline{V}$.

\item\label{theo5.6_2} The morphisms $\varphi:\Lambda_{(k)}\rightarrow\mathbb{R}$ nonnegative
on the $k$-Schur functions are parametrized by $\mathbb{R}_{\geq0}^{k}$.
\end{enumerate}
\end{theorem}

\begin{proof}
\eqref{theo5.6_1}: The set $A(\vec{r}_{0})$ is a connected subset by the previous lemma and
it is also finite by Proposition~\ref{Prop_Variety}. Therefore, the set
$A(\vec{r}_{0})$ is a singleton. In particular, $f(\vec{r})$ converges to some
$f(\vec{r}_{0})$ as $\vec{r}$ tends to $\vec{r}_{0}$, and $f$ can be extended
by continuity on $\mathbb{R}_{\geq0}^{k}$. As explained at the end of
\S~\ref{SubSec_Twoparam}, this suffices to conclude that $f$ is an
homeomorphism from $\overline{U}$ to $\overline{V}$.

\eqref{theo5.6_2}: By the first part of the theorem, it suffices to prove that any morphism
$\varphi:\Lambda_{(k)}\rightarrow\mathbb{R}$ nonnegative on the $k$-Schur
functions is in the closure of the set of positive morphisms. Let $\varphi
_{0}:\Lambda_{(k)}\rightarrow\mathbb{R}$ be a positive morphism (for example,
one can take the image of any element of $U$ by $f$). Then, for any
nonnegative morphism $\varphi:\Lambda_{(k)}\rightarrow\mathbb{R}$ and $0\leq
t\leq1$, one can define the convolution morphism $\varphi*_{t}\varphi_{0}$ by
the formula
\[
\varphi*_{t}\varphi_{0}(f)=((1-t)\cdot\varphi\otimes t\cdot\varphi_{0}%
)\Delta(f)
\]
for $f\in\Lambda_{(k)}$, where $\Delta$ is the coproduct on $\Lambda_{(k)}$.
Since $\Delta$ is an algebra morphism from $\Lambda_{(k)}$ to $\Lambda
_{(k)}\otimes\Lambda_{(k)}$, $\varphi*_{t}\varphi_{0}$ is indeed a morphism.
Let $\lambda\in\mathcal{B}_{k}$ and $s_{\lambda}^{(k)}$ the corresponding
$k$-Schur function. Then, by~\cite[Corollary~8.1]{LamLR},
\[
\Delta(s_{\lambda}^{(k)})=\sum_{\substack{\mu,\nu\in\mathcal{B}_{k}\\\vert
\mu\vert+\vert\nu\vert=\vert\lambda\vert}}C_{\mu,\nu}^{\lambda,(k)}s_{\mu
}^{(k)}\otimes s_{\nu}^{(k)}%
\]
with nonnegative coefficients $C_{\mu,\nu}^{\lambda,(k)}$. Since $\Delta$ is
the usual coproduct from the ring of symmetric functions, we moreover have
$C_{\lambda,\emptyset}^{\lambda,(k)}=C_{\emptyset,\lambda}^{\lambda,(k)}=1$,
so that
\[
\varphi*_{t}\varphi_{0}(s_{\lambda}^{(k)})=(1-t)\cdot\varphi(s_{\lambda}%
^{(k)})+t\cdot\varphi_{0}(s_{\lambda}^{(k)})+\sum_{\substack{\vert\mu
\vert+\vert\nu\vert=\vert\lambda\vert\\\mu,\nu\not =\emptyset}}C_{\mu,\nu
}^{\lambda,(k)}(1-t)^{\vert\mu\vert}\varphi(s_{\mu}^{(k)})t^{\vert\nu\vert}
\varphi_{0}(s_{\nu}^{(k)}).
\]
Since $\varphi_{0}(s_{\lambda}^{(k)})$ is positive and all terms in the above
sums are nonnegative, $\varphi*_{t}\varphi_{0}(s_{\lambda}^{(k)})$ is positive
for all $t>0$. Hence, $\varphi*_{t}\varphi_{0}$ is a positive morphism and
$\varphi*_{t}\varphi_{0}$ converges to $\varphi$ as $t$ goes to zero. Thus,
$\varphi$ is in the closure of the set of positive morphisms.
\end{proof}

\begin{thebibliography}{99}
\bibitem[5]{KR} Martin Kreuzer and Lorenzo Robbiano, \emph{Computational commutative algebra.2}, Springer-Verlag, Berlin,2005.
\bibitem[7]{LamLR} Thomas Lam, ``Schubert polynomials for the affine Grassmannian'', \emph{J. Amer. Math. Soc.}\textbf{21}(2008),no.1,259--281.
\bibitem[9]{LLMSSZ} Thomas Lam, Luc Lapointe, Jennifer Morse, Anne Schilling, Mark Shimozono and Mike Zabrocki, \emph{k-Schur functions and affine Schubert calculus}, Fields Institute Monographs33, Springer, New York; Fields Institute for Research in Mathematical Sciences, Toronto,ON,2014.
\bibitem[13]{LT} Cédric Lecouvey and Pierre Tarrago, ``Harmonic functions on multiplicative graphs and weight polytopes of representations'', To appear in Annales de l'Institut Fourier, \url{https://arxiv.org/abs/1609.00138}.
\bibitem[14]{Mac} Ian G. Macdonald, \emph{Symmetric functions and Hall polynomials}, second ed., Oxford Mathematical Monographs, The Clarendon Press, Oxford University Press, New York,1995.
\bibitem[15]{Mitt} Johannes Mittmann, \emph{Independence in algebraic complexity theory}, Ph.D.thesis, Universitäts- und Landesbibliothek Bonn,2013.
\bibitem[18]{Pan} Anurag Pandey, \emph{Algebraic independence: Criteria and structural results over diverse fields}, Master's thesis, Indian Institute of Technology Kanpur,2015.
\end{thebibliography}
\end{document}
